\documentclass[11pt,a4paper]{article}
\usepackage[margin=24mm]{geometry}
\usepackage[T1]{fontenc}
\usepackage[utf8]{inputenc}
\usepackage{lmodern}
\usepackage{amsmath,amssymb,amsthm,mathtools}
\usepackage{array,longtable,booktabs}
\usepackage{xurl}
\usepackage[unicode,hidelinks]{hyperref}
\usepackage{microtype}
\setlength{\parindent}{0pt}
\setlength{\parskip}{5pt}
\setlength{\emergencystretch}{3em}
\allowdisplaybreaks
\newcommand{\R}{\mathbb R}
\newcommand{\C}{\mathbb C}
\newcommand{\N}{\mathbb N}
\newcommand{\norm}[1]{\left\lVert#1\right\rVert}
\newcommand{\abs}[1]{\left|#1\right|}
\DeclareMathOperator{\tr}{tr}
\DeclareMathOperator{\diag}{diag}
\DeclareMathOperator{\rank}{rank}
\DeclareRobustCommand{\code}[1]{{\footnotesize\urlstyle{tt}\nolinkurl{#1}}}
\newtheorem{theorem}{Theorem}[section]
\newtheorem{lemma}[theorem]{Lemma}
\title{Open Matrix Analysis Toolkit Items\\
Finite and Operator Proof Routes}
\author{Mathematical derivations and audit for NLAlib}
\date{15 October 2026}
\begin{document}
\maketitle

\textbf{Scope.} The live
\href{https://research.chen.pw/NLAlib/}{NLAlib catalogue}
on 15 October 2026 has 176 matrix-toolkit cards: 107 proved and 69 open.
The user confirmed that none of the matrix-toolkit cards are deferred.
This standalone document covers all 69 open cards, including the exact
mathematical identities and operation counts attached to iterative methods.
It is intended for a separate Lean formalization effort; no Lean source or
catalogue status is changed here.

\textbf{Outcome and limitation.} Sixty-eight cards receive explicit
mathematical derivations, using shared finite cores wherever possible.
For Stahl's trace-exponential positive-measure theorem, the downstream
claims are proved and an equivalent finite-inequality interface is given,
but the general noncommutative positivity core remains an imported deep
theorem. This is identified explicitly in Section~\ref{route:trace-exponential-stahl}.
An explicit mathematical route does not mean its Lean implementation is
already easy: spectral-calculus, finite-dimensional topology, and integral
interfaces still require assembly. None of these catalogue entries is being
declared an unsolved research problem.

\textbf{Snapshot and reuse.} The site's embedded atlas was compared with
the repository atlas at commit
\href{https://github.com/tchen-research/NLAlib/tree/6886918346afb81ba77e71fbbe5daa67c93d2bae}
{\texttt{6886918346afb81ba77e71fbbe5daa67c93d2bae}} and the parsed data agree.
Reusable declarations were searched in the site's declaration index;
named NLAlib declarations were checked for \code{sorry=false}.
Specific Mathlib names are identified where their local source was checked.
These are reuse pointers, not a claim that new Lean applications were compiled.
The public atlas ids below give a stable lookup key even if a module is moved.

\textbf{Audit points.} The proofs retain repeated eigenvalues, rank-deficient
cases, and empty dimensions wherever the statements allow them. Spectral
norms and Frobenius norms are kept distinct; $*$ denotes an adjoint.
Real spectral assertions use complexification when nonreal eigenvalues
matter. In particular, real-only spectra would make the Sylvester
separation criterion false. The partial-pivoting sharpness example matches
the implementation's smallest-index tie rule. The polar convergence exponent
is $(q+1)^T$ for degree $2q+1$. No asymptotic cost statement includes
coefficient computation or floating-point analysis unless explicitly said.

\tableofcontents
\clearpage
\section*{Coverage ledger}
\addcontentsline{toc}{section}{Coverage ledger}
Every open card occurs once in the derivation body. ``Derived'' describes
the mathematics supplied here, not a completed formal proof.
\small
\begin{longtable}{@{}r p{0.60\textwidth} l l@{}}
\toprule
No. & Atlas id & Section & Result \\
\midrule
\endfirsthead
\toprule
No. & Atlas id & Section & Result \\
\midrule
\endhead
1 & \code{lu-growth-factor} & \ref{route:lu-growth-factor} & Derived \\
2 & \code{hessenberg-reduction} & \ref{route:hessenberg-reduction} & Derived \\
3 & \code{laplacian-resistance-leverage} & \ref{sec:ns-resistance} & Derived \\
4 & \code{operator-monotone-modulus-continuity} & \ref{sec:operator-monotone-modulus-continuity} & Derived \\
5 & \code{ordered-schatten-trace-deficit} & \ref{sec:ordered-schatten-trace-deficit} & Derived \\
6 & \code{tikhonov-spectral-filter} & \ref{sec:eq-tikhonov} & Derived \\
7 & \code{truncated-svd-regularization} & \ref{sec:eq-tsvd} & Derived \\
8 & \code{landweber-spectral-filter} & \ref{sec:eq-landweber} & Derived \\
9 & \code{displacement-rank-def} & \ref{route:displacement-rank-def} & Derived \\
10 & \code{displacement-polynomial-rank} & \ref{route:displacement-polynomial-rank} & Derived \\
11 & \code{zolotarev-number-def} & \ref{route:zolotarev-number-def} & Derived \\
12 & \code{displacement-singular-value-decay} & \ref{route:displacement-singular-value-decay} & Derived \\
13 & \code{positive-map-kadison} & \ref{sec:positive-map-kadison} & Derived \\
14 & \code{positive-map-inverse-jensen} & \ref{sec:positive-map-inverse-jensen} & Derived \\
15 & \code{positive-map-norm-at-identity} & \ref{sec:positive-map-norm-at-identity} & Derived \\
16 & \code{matrix-divided-differences-def} & \ref{sec:matrix-divided-differences-def} & Derived \\
17 & \code{matrix-function-frechet-divided-difference} & \ref{sec:matrix-function-frechet-divided-difference} & Derived \\
18 & \code{matrix-monotone-loewner-criterion} & \ref{sec:matrix-monotone-loewner-criterion} & Derived \\
19 & \code{matrix-exponential-frechet-integral} & \ref{sec:matrix-exponential-frechet-integral} & Derived \\
20 & \code{matrix-geometric-mean-def} & \ref{sec:matrix-geometric-mean-def} & Derived \\
21 & \code{matrix-geometric-mean-block-maximum} & \ref{sec:matrix-geometric-mean-block-maximum} & Derived \\
22 & \code{schatten-even-moment-smoothness} & \ref{sec:schatten-even-moment-smoothness} & Derived \\
23 & \code{matrix-sign-spectral-projectors} & \ref{sec:eq-sign-projectors} & Derived \\
24 & \code{matrix-sign-newton} & \ref{sec:eq-sign-newton} & Derived \\
25 & \code{matrix-sign-sylvester-block} & \ref{sec:eq-sign-sylvester} & Derived \\
26 & \code{matrix-inverse-newton-schulz} & \ref{sec:eq-inverse-schulz} & Derived \\
27 & \code{polar-newton-iteration} & \ref{sec:eq-polar-newton} & Derived \\
28 & \code{polar-newton-schulz-iteration} & \ref{sec:eq-polar-schulz} & Derived \\
29 & \code{nearest-psd-frobenius} & \ref{sec:eq-nearest-psd} & Derived \\
30 & \code{pseudospectrum-perturbation-equivalence} & \ref{sec:eq-pseudospectrum} & Derived \\
31 & \code{matrix-exponential-logarithmic-norm} & \ref{sec:eq-lognorm} & Derived \\
32 & \code{gsvd-normalized-pair} & \ref{sec:eq-gsvd} & Derived \\
33 & \code{general-form-tikhonov-gsvd} & \ref{sec:eq-general-tikhonov} & Derived \\
34 & \code{total-least-squares-svd} & \ref{sec:eq-tls} & Derived \\
35 & \code{generalized-schur-qz} & \ref{sec:eq-qz} & Derived \\
36 & \code{kkt-reduced-hessian} & \ref{sec:eq-kkt} & Derived \\
37 & \code{lyapunov-stability-gramian} & \ref{route:lyapunov-stability-gramian} & Derived \\
38 & \code{sylvester-separation-residual-conditioning} & \ref{sec:eq-separation} & Derived \\
39 & \code{adi-sylvester-rational-error} & \ref{route:adi-sylvester-rational-error} & Derived \\
40 & \code{lyapunov-adi-residual-factor} & \ref{route:lyapunov-adi-residual-factor} & Derived \\
41 & \code{stein-gramian-smith-doubling} & \ref{route:stein-gramian-smith-doubling} & Derived \\
42 & \code{matrix-function-block-divided-differences} & \ref{sec:matrix-function-block-divided-differences} & Derived \\
43 & \code{trace-exponential-stahl} & \ref{route:trace-exponential-stahl} & Partial \\
44 & \code{nonnegative-perron-eigenpair} & \ref{sec:ns-perron} & Derived \\
45 & \code{irreducible-perron-frobenius} & \ref{sec:ns-perron} & Derived \\
46 & \code{nonnegative-spectral-radius-comparison} & \ref{sec:ns-perron} & Derived \\
47 & \code{collatz-wielandt-certificates} & \ref{sec:ns-perron} & Derived \\
48 & \code{primitive-perron-power-limit} & \ref{sec:ns-perron} & Derived \\
49 & \code{nonsingular-m-matrix-characterizations} & \ref{sec:ns-mmatrix} & Derived \\
50 & \code{m-matrix-schur-complement} & \ref{sec:ns-mmatrix} & Derived \\
51 & \code{circulant-fourier-diagonalization} & \ref{sec:ns-fourier} & Derived \\
52 & \code{circulant-fourier-pseudoinverse} & \ref{sec:ns-fourier} & Derived \\
53 & \code{radix-two-fft-realization} & \ref{sec:ns-fourier} & Derived \\
54 & \code{toeplitz-circulant-fft-product} & \ref{sec:ns-fourier} & Derived \\
55 & \code{dmd-linear-consistency} & \ref{route:dmd-linear-consistency} & Derived \\
56 & \code{dmd-reduced-eigenpair-lifting} & \ref{route:dmd-reduced-eigenpair-lifting} & Derived \\
57 & \code{hankel-svd-era-realization} & \ref{route:hankel-svd-era-realization} & Derived \\
58 & \code{gramian-svd-balancing} & \ref{route:gramian-svd-balancing} & Derived \\
59 & \code{balanced-truncation-energy-bound} & \ref{route:balanced-truncation-energy-bound} & Derived \\
60 & \code{deim-greedy-interpolation} & \ref{route:deim-greedy-interpolation} & Derived \\
61 & \code{hodlr-storage-matvec} & \ref{sec:ns-hierarchical} & Derived \\
62 & \code{hodlr-levelwise-compression-error} & \ref{sec:ns-hierarchical} & Derived \\
63 & \code{hss-nested-bases-matvec} & \ref{sec:ns-hierarchical} & Derived \\
64 & \code{block-separable-inverse-reduction} & \ref{sec:ns-hierarchical} & Derived \\
65 & \code{hss-inverse-rank-preservation} & \ref{sec:ns-hierarchical} & Derived \\
66 & \code{polar-minimax-polynomial-iteration} & \ref{route:polar-minimax-polynomial-iteration} & Derived \\
67 & \code{normal-matrix-mixed-chebyshev-moments} & \ref{sec:ns-mixedmoments} & Derived \\
68 & \code{matrix-k-transform-stability} & \ref{sec:matrix-k-transform-stability} & Derived \\
69 & \code{coefficient-block-partial-transpose} & \ref{sec:coefficient-block-partial-transpose} & Derived \\
\bottomrule
\end{longtable}
\normalsize
\clearpage
\section{Two partially completed factorization cards}
\subsection{Sharp partial-pivoting growth}
\label{route:lu-growth-factor}
\textbf{Card:} \code{lu-growth-factor}.
The upper bound is already available as
\code{NLAlib.abs_luPartialPivot_upper_le_two_pow}; only attainment needs a new argument.
For $n\ge1$, use zero-based indices and define
\[
 W_{ij}=\begin{cases}1&j=n-1,\\1&i=j<n-1,\\-1&i>j,\ j<n-1,\\0&i<j<n-1.\end{cases}
\]
At the start of elimination step $k<n-1$, the trailing leading block still
has diagonal $1$, strict lower triangle $-1$, and strict upper triangle $0$;
its last column is constantly $2^k$. Thus the current diagonal is a largest
modulus entry of the pivot column. Choosing it gives multipliers $-1$ and
doubles every remaining last-column entry. Induction gives
\[
 U_{ii}=1\ (i<n-1),\qquad U_{i,n-1}=2^i,
 \qquad U_{ij}=0\ (i<j<n-1).
\]
Consequently $\max|W_{ij}|=1$ and $\max|U_{ij}|=2^{n-1}$.
This is the actual tie convention of the inspected implementation:
\code{absArgmax} chooses the smallest index of largest modulus, so the
canonical \code{luPartialPivot} performs no swaps on this example.
For formalization, prove a stronger induction with last column constantly
$c$: one Schur step replaces $c$ by $2c$ and reduces the dimension by one.
The case $n=1$ is $[1]$; for $n=0$ retain the entrywise bound and do not
define a growth ratio by division by the maximum of an empty family.

\subsection{Hessenberg columns are the Arnoldi basis after signs are fixed}
\label{route:hessenberg-reduction}
\textbf{Card:} \code{hessenberg-reduction}.
Suppose $AQ=QH$, $Q$ is real orthogonal, and $H$ is upper Hessenberg.
Write its columns as $q_0,\ldots,q_{n-1}$. For $j<n-1$,
\[
 Aq_j=\sum_{i=0}^{j+1}h_{ij}q_i.
\]
If $0\le s\le n-1$ and $h_{j+1,j}\ne0$ up to step $s-1$, induction gives
\[
 A^jq_0=\left(\prod_{i=0}^{j-1}h_{i+1,i}\right)q_j
       +\operatorname{span}(q_0,\ldots,q_{j-1}),\qquad j\le s.
\]
The triangular change-of-basis matrix has nonzero diagonal; hence
$\operatorname{span}(q_0,Aq_0,\ldots,A^jq_0)
=\operatorname{span}(q_0,\ldots,q_j)$.
Set $d_0=1$ and $d_{j+1}=d_j\operatorname{sign}(h_{j+1,j})$ for
$j<s$, and extend the remaining $d_j$ arbitrarily in $\{-1,1\}$.
With $D=\diag(d_j)$, the columns of $QD$ have the same first column and
$DHD$ has positive subdiagonal on this prefix. Their Arnoldi residual is exactly
$|h_{j+1,j}|(d_{j+1}q_{j+1})$, so normalization proves equality with the
Arnoldi recursion, one column at a time. There is no need to formalize
uniqueness of a whole Krylov factorization first.
If a subdiagonal vanishes, this identifies only the prefix before breakdown;
the prefix span is then invariant and later columns are an orthogonal
completion, not further nonzero Arnoldi residuals. Dimension zero is vacuous.
Reuse \code{NLAlib.IsHessenbergForm.exists_sign} and
\code{NLAlib.isHessenbergForm_hessenbergReduction} for the already proved
reduction and implicit-$Q$ portions.

\section{Displacement rank from one bivariate polynomial identity}
\subsection{Generators and the rank convention}
\label{route:displacement-rank-def}
\textbf{Card:} \code{displacement-rank-def}.
For $A\in\C^{m\times m}$, $B\in\C^{n\times n}$ and $X\in\C^{m\times n}$,
put $\Delta(X)=AX-XB$. The condition $\rank\Delta(X)\le\nu$ is equivalent
to $\Delta(X)=MN^*$ with $M\in\C^{m\times\nu}$ and
$N\in\C^{n\times\nu}$. One implication is the product rank inequality.
For the other, choose a basis of the column space of $\Delta(X)$,
write every column in that basis, and pad the two factors by zero columns.
This also handles $\nu=0$ and empty row or column types.

\subsection{The degree bound is $k\nu$, not $k^2\nu$}
\label{route:displacement-polynomial-rank}
\textbf{Card:} \code{displacement-polynomial-rank}.
If $\deg p,\deg q\le k$ and $k\ge1$, then
\[
 p(z)q(w)-q(z)p(w)=(z-w)\sum_{i,j=0}^{k-1}c_{ij}z^iw^j. \tag{D1}
\]
Divisibility follows by setting $z=w$. The degree in $z$ of the quotient
is at most $k-1$, and antisymmetry of the numerator makes the quotient
symmetric, proving the same bound in $w$. Substitution into the bimodule
of rectangular matrices gives
\[
 p(A)Xq(B)-q(A)Xp(B)
 =\sum_{i,j<k}c_{ij}A^i(AX-XB)B^j.
\]
If $AX-XB=MN^*$, every column on the right belongs to
$\operatorname{range}[M,AM,\ldots,A^{k-1}M]$, a space of dimension at most
$k\nu$. This is one column-space inclusion, rather than a rank estimate
for $k^2$ separate summands. When $k=0$ the expression is zero.
If $q(A)$ and $p(B)$ are invertible, multiplication by their inverses gives
\[
 \rank\bigl(X-q(A)^{-1}p(A)Xq(B)p(B)^{-1}\bigr)\le k\nu. \tag{D2}
\]
No inverse of $p(A)$ or $q(B)$ is required.

\subsection{Zolotarev numbers without minimizers or pole bookkeeping}
\label{route:zolotarev-number-def}
\textbf{Card:} \code{zolotarev-number-def}.
For nonempty disjoint compact sets $E,F\subset\C$, define
\[
 Z_k(E,F)=\inf_{(p,q)\in\mathcal A_k}
  \max_{z\in E}\abs{p(z)/q(z)}\max_{w\in F}\abs{q(w)/p(w)},
\]
where $\mathcal A_k$ consists of polynomials of degrees at most $k$,
with $q$ nowhere zero on $E$ and $p$ nowhere zero on $F$.
This pair formulation permits poles on $F$ and avoids an extended-valued
rational function. Constants $(1,1)$ show nonemptiness and $Z_k\le1$;
all competitors are nonnegative, so $0\le Z_k\le1$ and $Z_0=1$.
Inclusion of degree classes gives $Z_{k+1}\le Z_k$.
Enlarging either set restricts admissibility and increases each surviving
supremum, so it can only increase $Z_k$.
Swapping $(E,F)$ and $(p,q)$ proves symmetry.
Multiplying two admissible pairs proves
\[
 Z_{k+\ell}(E,F)\le Z_k(E,F)Z_\ell(E,F).
\]
For precision, choose competitors within $\epsilon$ of each infimum,
multiply, and let $\epsilon\downarrow0$. This works when either infimum
is zero and makes no attainment claim.

\subsection{Singular-value decay by a rank correction}
\label{route:displacement-singular-value-decay}
\textbf{Card:} \code{displacement-singular-value-decay}.
Assume $A,B$ are normal, their spectra lie in $E,F$, and
$\rank(AX-XB)\le\nu$. Use descending singular values indexed from $1$,
extended by zero beyond the matrix dimensions. For an admissible $(p,q)$,
put $L=q(A)^{-1}p(A)$, $R=q(B)p(B)^{-1}$, and $D=X-LXR$.
By (D2), $\rank D\le k\nu$. Let $X_{j-1}$ be the truncated SVD of rank
at most $j-1$. The matrix $D+LX_{j-1}R$ has rank at most $k\nu+j-1$, so
\[
 \sigma_{j+k\nu}(X)
 \le\norm{L(X-X_{j-1})R}_2
 \le\norm L_2\norm R_2\sigma_j(X).
\]
Normal spectral calculus bounds the two norms by the maxima in the
definition of $Z_k$. Taking the infimum, without dividing by
$\sigma_j(X)$, proves
\[
 \boxed{\sigma_{j+k\nu}(X)\le Z_k(E,F)\sigma_j(X).}
\]
If $\sigma_j(X)=0$, the displayed inequality for any one competitor already
gives the result. Reuse
\code{NLAlib.singularValues_le_specNorm_sub_of_rank_le} for the low-rank
approximation inequality. The algebraic core follows the displacement
viewpoint of Beckermann--Townsend, but needs no contour calculus or explicit
elliptic-function formula for $Z_k$.

\section{Continuous and discrete Gramians and ADI}
\subsection{Continuous Lyapunov stability and controllability}
\label{route:lyapunov-stability-gramian}
\textbf{Card:} \code{lyapunov-stability-gramian}.
Suppose $A$ is Hurwitz and $Q\succeq0$. A finite triangular or Jordan
decomposition gives $\norm{e^{tA}}\le C(1+t^{n-1})e^{-ct}$ for some
$c>0$. Indeed each Jordan block contributes
$e^{t\lambda}\sum_{j<s}t^jN^j/j!$; alternatively triangular variation of
constants gives the same bound without selecting Jordan chains.
Thus the entrywise integral
\[
 X=\int_0^\infty e^{tA}Qe^{tA^*}\,dt \tag{G1}
\]
converges absolutely. Differentiation and the fundamental theorem of
calculus give $AX+XA^*=-Q$. Each integrand is PSD. If $Q\succ0$, its
quadratic form at $t=0$ is positive for every nonzero vector, hence $X\succ0$.
If $AY+YA^*=0$, then $e^{tA}Ye^{tA^*}$ is constant and tends to zero;
therefore $Y=0$, proving uniqueness without a Kronecker inverse.

Conversely, if $X\succ0$ and $AX+XA^*\prec0$, choose an eigenvector
$A^*v=\overline\lambda v$. Its quadratic form is
$2\operatorname{Re}\lambda\,v^*Xv<0$, so every eigenvalue has negative
real part. This works for real matrices by complexification.
For $Q=BB^*$, (G1) gives
\[
 v^*Xv=\int_0^\infty\norm{B^*e^{tA^*}v}^2dt.
\]
A continuous nonnegative integrand has zero integral exactly when it is
identically zero. Differentiating at zero and using Cayley--Hamilton gives
\[
 \ker X=\bigcap_{j=0}^{n-1}\ker(B^*(A^*)^j),\qquad
 \rank X=\rank[B,AB,\ldots,A^{n-1}B]. \tag{G2}
\]
For the reverse implication in (G2), the first $n$ vanishing powers imply
all powers vanish by the characteristic recurrence, hence the exponential
series vanishes. A low-rank $Q$ does not imply a low-rank Gramian.

\subsection{General Sylvester ADI: exact error, residual, and rank}
\label{route:adi-sylvester-rational-error}
\textbf{Card:} \code{adi-sylvester-rational-error}.
Let $AX-XB=C$. At step $j$ assume only $A-\beta_jI$ and $B-\alpha_jI$
are invertible, and define
\[
 (A-\beta_jI)X_{j-1/2}=X_{j-1}(B-\beta_jI)+C,
 \quad X_j(B-\alpha_jI)=(A-\alpha_jI)X_{j-1/2}-C.
\]
Subtract the same equations for the exact $X$ to obtain, with $E_j=X-X_j$,
\[
 E_j=(A-\alpha_jI)(A-\beta_jI)^{-1}E_{j-1}
                (B-\beta_jI)(B-\alpha_jI)^{-1}.
\]
All factors on the left are functions of $A$, and all factors on the right
are functions of $B$. Hence $E_j=F_jE_0G_j$, with the corresponding products
$F_j,G_j$. Submultiplicativity gives
$\norm{E_j}\le\norm{F_j}_2\norm{G_j}_2\norm{E_0}$ for either
the spectral or Frobenius norm on the error. The residual
$R_j=C-AX_j+X_jB=AE_j-E_jB$ obeys
$R_j=F_jR_0G_j$. Direct subtraction also gives the useful update
\[
 X_j-X_{j-1}=(\beta_j-\alpha_j)(A-\beta_jI)^{-1}R_{j-1}
                                         (B-\alpha_jI)^{-1}. \tag{G3}
\]
Thus $\rank R_j\le\rank R_0$, and if $X_0=0$, $\rank C\le\nu$, then
$\rank X_j\le j\nu$. For normal $A,B$, spectral calculus turns the
norms of $F_j,G_j$ into maxima of scalar products on their spectra.
These are identities and bounds for arbitrary admissible shifts, not a
convergence assertion. Equal shifts are allowed and make (G3) zero;
the unused numerator factors need not be invertible.

\subsection{Positive Lyapunov ADI factors}
\label{route:lyapunov-adi-residual-factor}
\textbf{Card:} \code{lyapunov-adi-residual-factor}.
Let $AX+XA^*+BB^*=0$, with $A$ Hurwitz, and choose $p_j>0$.
Set $W_0=B$, $X_0=0$ and
\[
 V_j=(A-p_jI)^{-1}W_{j-1},\quad
 W_j=W_{j-1}+2p_jV_j,\quad
 X_j=X_{j-1}+2p_jV_jV_j^*.
\]
Since $W_{j-1}=(A-p_jI)V_j$, expanding $W_jW_j^*$ shows
\[
 AX_j+X_jA^*+BB^*=W_jW_j^*. \tag{G4}
\]
For $F_j=\prod_{i=1}^j(A+p_iI)(A-p_iI)^{-1}$, either the same expansion
or the Sylvester error identity proves
$X-X_j=F_jXF_j^*$. Therefore
\[
 0\preceq X_{j-1}\preceq X_j\preceq X,\qquad
 X_j=Z_jZ_j^*,\quad Z_j=[\sqrt{2p_1}V_1,\ldots,\sqrt{2p_j}V_j].
\]
Ranks are at most $j\rank B$ and hence at most $j$ times the number of
columns of $B$. The residual norms can be computed with small factors:
$\norm{W_jW_j^*}_2=\norm{W_j}_2^2$ and
$\norm{W_jW_j^*}_F=\norm{W_j^*W_j}_F$.
For a fixed positive shift $p$, the eigenvalues of the Cayley matrix are
$(\lambda+p)/(\lambda-p)$ and have modulus less than one. Finite
triangular power decay implies $F_j\to0$, even for nonnormal $A$.
No convergence for an unrestricted varying sequence of positive shifts
is asserted.

\subsection{Stein equations and Smith doubling}
\label{route:stein-gramian-smith-doubling}
\textbf{Card:} \code{stein-gramian-smith-doubling}.
For $\rho(A)<1$ and $Q\succeq0$, set
\[
 X_N=\sum_{j=0}^{N-1}A^jQ(A^*)^j,\qquad
 X=\sum_{j=0}^{\infty}A^jQ(A^*)^j.
\]
Geometric decay with a polynomial prefactor proves absolute convergence.
Index shifting gives
\[
 X-AXA^*=Q,\quad X-X_N=A^NX(A^*)^N,\quad
 Q-X_N+AX_NA^*=A^NQ(A^*)^N. \tag{G5}
\]
A homogeneous solution satisfies $Y=A^NY(A^*)^N$ and is zero. Positivity
and monotonicity follow term by term. Starting with $A_0=A,G_0=Q$, the rules
\[
 A_{k+1}=A_k^2,\qquad G_{k+1}=G_k+A_kG_kA_k^*
\]
give $A_k=A^{2^k}$ and $G_k=X_{2^k}$ by splitting a finite sum in half.
For either spectral or Frobenius norm, $\norm A_2\le q<1$ yields
$\norm{X-X_N}\le q^{2N}\norm Q/(1-q^2)$.
If $G_k=Z_kZ_k^*$, take $Z_{k+1}=[Z_k,A_kZ_k]$; its column count doubles
unless compression is added, so this identity alone is not a complexity
claim. Finally $\rho(A)<1$ is equivalent to existence of $X\succ0$ with
$X-AXA^*\succ0$: the forward direction uses $Q=I$, and the reverse direction
tests an eigenvector of $A^*$, yielding $(1-|\lambda|^2)v^*Xv>0$.
The existing declarations
\code{NLAlib.exists_prefix_norm_matrix_pow_le_spectralRadius_add} and
\code{NLAlib.hasSum_matrix_pow_inv_of_spectralRadius_lt_one}
provide the relevant finite-dimensional power-decay/Neumann interfaces.


\section{Matrix functions: three finite cores}
\label{sec:functions-finite-cores}

The following fifteen entries reduce chiefly to spectral projectors, finite
polynomial identities, and quadratic forms. All norms without a subscript are
spectral norms. A Schatten norm uses the unnormalised trace. The adjoint is
denoted by $*$, over either $\R$ or $\C$ as appropriate. Empty matrix spaces
make all norm identities below trivial; when a proof chooses an eigenvalue or
a unit vector, first separate that case. Strictly positive matrix means positive
definite, rather than a matrix with positive entries.

\subsection{Operator-monotone modulus of continuity}
\label{sec:operator-monotone-modulus-continuity}
\textbf{Atlas:} \code{operator-monotone-modulus-continuity}.
Let $f:[0,\infty)\to[0,\infty)$ be continuous, operator monotone in every
dimension, and satisfy $f(0)=0$. Then, for $A,B\succeq0$,
\[
                 \norm{f(A)-f(B)}\le f(\norm{A-B}).                 \tag{F1}
\]
Here is a small scalar-concavity lemma that avoids importing an integral
representation of operator-monotone functions. Fix $x,y\ge0$, $0\le s\le1$,
and conjugate $\diag(x,y)$ by a real rotation whose first row has squared
entries $s,1-s$. Write the resulting PSD matrix as
\[
 D=\begin{bmatrix}a&b\\b&c\end{bmatrix},\qquad a=sx+(1-s)y.
\]
For every $\varepsilon>0$, completing the square gives
\[
 D\preceq\diag(a+\varepsilon,c+b^2/\varepsilon).
\]
Operator monotonicity, followed by the first diagonal entry, yields
$sf(x)+(1-s)f(y)\le f(a+\varepsilon)$. Letting $\varepsilon\downarrow0$
proves scalar concavity, including $a=0$. Apply concavity twice on
$[0,x+y]$ to obtain
\[
 f(x)\ge\frac{x}{x+y}f(x+y),\qquad
 f(y)\ge\frac{y}{x+y}f(x+y),
\]
when $x+y>0$; the zero case is immediate. Hence
$f(x+y)\le f(x)+f(y)$.

Put $t=\norm{A-B}$. Hermitian norm/order equivalence gives
$A\preceq B+tI$ and $B\preceq A+tI$. Consequently
\[
 f(A)\preceq f(B+tI)\preceq f(B)+f(t)I.
\]
The second step is the scalar inequality just proved, applied in an eigenbasis
of $B$; it does not assert operator subadditivity for unrelated matrices.
Interchange $A,B$ and apply norm/order equivalence again to prove (F1).
For $t=0$ this directly says $f(A)=f(B)$. The proof needs only the assumed
matrix monotonicity, two-dimensional scalar concavity, and simultaneous
functional calculus for $B,B+tI$.

\subsection{Schatten deficits from one padded-spectrum majorization}
\label{sec:ordered-schatten-trace-deficit}
\textbf{Atlas:} \code{ordered-schatten-trace-deficit}.
For $X,Y\succeq0$ of size $n$,
\[
 (\lambda(X),\lambda(Y))\prec(\lambda(X+Y),0_n).                 \tag{F2}
\]
To prove the top-$k$ condition when $k\le n$, select any $k$ coordinates
from the two spectra. Let $W$ contain their corresponding eigenvectors,
and enlarge it to dimension $k$. If $P$ is its orthogonal projector, PSD
gives
\[
 \text{selected sum}\le\tr(PX)+\tr(PY)
 =\tr(P(X+Y))\le\sum_{j=1}^k\lambda_j(X+Y).
\]
For example, the first inequality follows by expanding $\tr(PX)$ in an
orthonormal eigenbasis of $X$: the selected vectors have projection length
one, and every remaining contribution is nonnegative. For $n\le k\le2n$
the right side of the majorization inequality is the total trace, and every
selected sum on the left is at most that trace. The total sums are equal.
Thus (F2) includes repeated and zero eigenvalues without a limiting argument.

Convexity of $u\mapsto u^p$ on $[0,\infty)$, $p\ge1$, now gives
\[
 \tr(X^p)+\tr(Y^p)\le\tr((X+Y)^p).
\]
With $X=A-B$, $Y=B$, this is precisely
\[
 \norm{A-B}_{S_p}^{p}\le\norm A_{S_p}^{p}-\norm B_{S_p}^{p}
 \quad(A\succeq B\succeq0).
\]
For $p=1$, trace linearity gives equality. No use is made of operator
monotonicity of $u^{p-1}$.

\paragraph{Verified reuse.}
\code{NLAlib.trace_compression_le_eigenvaluePrefix} is already polymorphic
over \code{RCLike}; its conclusion is
$\operatorname{Re}\tr(Q^*AQ)\le\mathrm{eigenvaluePrefix}(A,r)$ under
$A=A^*$ and $Q^*Q=I_r$.
\code{NLAlib.WeaklyMajorized.sum_rpow_le} takes nonnegative vectors,
weak majorization, and $1\le p$, and returns their power-sum inequality.
It therefore applies directly to the two nonnegative vectors in (F2).

\subsection{One spectral-projector lemma for three positive-map results}
\label{sec:positive-map-kadison}
\textbf{Atlas:} \code{positive-map-kadison},
\code{positive-map-inverse-jensen}, \code{positive-map-norm-at-identity}.
Let $\Phi:M_n(\C)\to M_k(\C)$ be positive and complex-linear. It preserves
adjoints: decompose a Hermitian matrix into its positive and negative spectral
parts to see that its image is Hermitian, and then write a general matrix as
$H+iK$.

The common finite lemma is the following. If $E_j\succeq0$ and $u_j,v_j\in\C$,
then
\[
 \sum_j\begin{bmatrix}|u_j|^2E_j&u_j\overline{v_j}E_j\\
 \overline{u_j}v_jE_j&|v_j|^2E_j\end{bmatrix}\succeq0.             \tag{F3}
\]
Indeed its quadratic form at $(x,y)$ is
$\sum_j(\overline{u_j}x+\overline{v_j}y)^*
 E_j(\overline{u_j}x+\overline{v_j}y)$. This proof is a finite sum and
requires no assertion that $\Phi$ is completely positive.

\paragraph{Kadison.}
Suppose $\Phi(I)=I$ and $A=\sum_j\lambda_jP_j$ is Hermitian. Use
$E_j=\Phi(P_j)$, $u_j=1$, $v_j=\lambda_j$ in (F3). The resulting block is
\[
                 \begin{bmatrix}I&\Phi(A)\\\Phi(A)&\Phi(A^2)\end{bmatrix}
                 \succeq0.
\]
Its quadratic form at $(-\Phi(A)y,y)$ gives
$\Phi(A^2)-\Phi(A)^2\succeq0$. This is the Hermitian-input inequality;
it does not imply the stronger Schwarz inequality on arbitrary inputs.

\paragraph{Inverse Jensen.}
\label{sec:positive-map-inverse-jensen}
For $A\succ0$, choose $a>0$ with $A\succeq aI$. Then
$\Phi(A)\succeq aI\succ0$. In (F3) use
$u_j=\sqrt{\lambda_j}$ and $v_j=1/\sqrt{\lambda_j}$. This gives
\[
                 \begin{bmatrix}\Phi(A)&I\\I&\Phi(A^{-1})\end{bmatrix}
                 \succeq0.
\]
Testing at $(-\Phi(A)^{-1}y,y)$ proves
$\Phi(A)^{-1}\preceq\Phi(A^{-1})$.

\paragraph{Norm at the identity, without regularization.}
\label{sec:positive-map-norm-at-identity}
Unitality is no longer assumed. Set $S=\Phi(I)\succeq0$. For a unitary
$U=\sum_j z_jP_j$, $|z_j|=1$, apply (F3) with $u_j=z_j$, $v_j=1$:
\[
                 \begin{bmatrix}S&\Phi(U)\\\Phi(U)^*&S\end{bmatrix}
                 \succeq0.                                      \tag{F4}
\]
For any PSD block $\left[\begin{smallmatrix}S&D\\D^*&S\end{smallmatrix}\right]$,
testing its quadratic form on $(tx,y)$ with arbitrary complex $t$ gives
\[
 |x^*Dy|^2\le(x^*Sx)(y^*Sy)\le\norm S^2\norm x^2\norm y^2.
\]
The assertion remains true if either quadratic coefficient is zero, since
then nonnegativity for all $t$ forces the cross term to vanish. Hence
$\norm{\Phi(U)}\le\norm S$, even if $S$ is singular or zero.

Every complex square contraction $A=V\diag(s_j)W^*$ is the average of
two unitaries, explicitly
\[
 U_\pm=V\diag(s_j\pm i\sqrt{1-s_j^2})W^*,\qquad A=(U_++U_-)/2.
\]
Thus $\norm{\Phi(A)}\le\norm S$. Since $\norm I=1$ when $n\ge1$,
evaluation at $I$ gives the reverse bound and proves
$\norm\Phi_{2\to2}=\norm{\Phi(I)}$. There is no inverse, support
compression, or $\varepsilon\tr(\cdot)I$ limit in this proof.

\paragraph{Verified reuse and small bridges.}
Use Mathlib's \code{PositiveLinearMap}, which extends a linear map and an
order homomorphism. Its \code{PositiveLinearMap.norm_apply_le_of_nonneg}
already gives $\norm{\Phi(X)}\le\norm{\Phi(I)}\norm X$ for $X\succeq0$;
it does not cover the arbitrary-input conclusion above.
\code{NLAlib.posSemidef_fromBlocks_iff} is proved over \code{RCLike}
and provides the generalized Schur criterion. The direct quadratic-form
tests above are shorter here than invoking that generalized criterion.
The indexed SVD and spectral theorem supply the finite decompositions.

\subsection{Divided differences, including coincident nodes}
\label{sec:matrix-divided-differences-def}
\textbf{Atlas:} \code{matrix-divided-differences-def}.
For $f\in C^1(I)$ on an open interval, define
\[
 f^{[1]}(x,y)=\begin{cases}(f(x)-f(y))/(x-y),&x\ne y,\\f'(x),&x=y.
 \end{cases}\qquad L_f(\lambda)_{ij}=f^{[1]}(\lambda_i,\lambda_j).
\]
The fundamental theorem of calculus gives the single formula
\[
 f^{[1]}(x,y)=\int_0^1f'((1-t)y+tx)\,dt.                       \tag{F5}
\]
It proves symmetry by $t\mapsto1-t$, joint continuity by uniform continuity
of $f'$ on a compact interval around any fixed pair, and the bound
$|f^{[1]}(x,y)|\le\sup_{[\min(x,y),\max(x,y)]}|f'|$.
The ordered endpoints in this supremum also cover $x>y$; at $x=y$ the
interval is a singleton. The diagonal value must remain $f'(x)$.
Mathlib's \code{deriv} supplies the total derivative convention; outside the
stated differentiability domain no mathematical property is asserted.

\paragraph{Verified reuse.}
For polynomials use \code{Polynomial.dividedDifference}.
\code{Polynomial.eval_dividedDifference_of_ne} gives the quotient at
unequal nodes, \code{Polynomial.eval_dividedDifference_self} gives the
polynomial derivative at the repeated node, and
\code{Polynomial.eval_dividedDifference_comm} gives symmetry. These are
already proved in the public declaration index.

\subsection{Daleckii--Krein by finite Hermite interpolation at one matrix}
\label{sec:matrix-function-frechet-divided-difference}
\textbf{Atlas:} \code{matrix-function-frechet-divided-difference}.
This route avoids approximation of $f'$ by polynomials and avoids any theorem
about uniform convergence of a sequence of derivatives.

Let $A=A^*$ have distinct eigenvalues $a_1,\ldots,a_s$ in $I$. Define
\[
 \ell_j(x)=\prod_{h\ne j}\frac{x-a_h}{a_j-a_h},\qquad
 p(x)=\sum_{j=1}^s\left[
 f(a_j)(1-2\ell_j'(a_j)(x-a_j))+f'(a_j)(x-a_j)
 \right]\ell_j(x)^2.                                           \tag{F6}
\]
Direct evaluation and differentiation show
$p(a_j)=f(a_j)$ and $p'(a_j)=f'(a_j)$. The products are over distinct
eigenvalues; repeated eigenvalues of $A$ create no denominators. Put $g=f-p$.
For every $\varepsilon>0$, differentiability at these finitely many points
gives a common $\delta>0$ such that
\[
 |g(x)|\le\varepsilon|x-a_j|\quad\text{if }|x-a_j|<\delta.       \tag{F7}
\]
Choose $\delta$ also small enough that all these neighbourhoods lie in $I$.

For a Hermitian perturbation $H$, every eigenvalue $\mu$ of $A+H$ lies
within $\norm H$ of some eigenvalue of $A$. A particularly small proof is
available over both fields: for a unit $\mu$-eigenvector $v$,
\[
 \norm{(A-\mu I)v}=\norm{Hv}\le\norm H;
\]
expanding $v$ in an orthonormal eigenbasis of $A$ shows that this norm is
at least $\min_j|a_j-\mu|$. Thus (F7) gives, for $\norm H<\delta$,
\[
 \norm{g(A+H)}=\max_{\mu\in\operatorname{spec}(A+H)}|g(\mu)|
 \le\varepsilon\norm H.                                      \tag{F8}
\]
Also $g(A)=0$. Equation (F8) is a full Fr\'echet remainder estimate on the
real vector space of Hermitian matrices, uniformly over all directions $H$.

For the polynomial $p(x)=\sum_r c_rx^r$, expanding products gives
\[
 Dp(A)[H]=\sum_{r\ge1}c_r\sum_{q=0}^{r-1}A^qHA^{r-1-q};        \tag{F9}
\]
the remaining terms contain at least two factors $H$ and have norm
$O(\norm H^2)$, since the sum is finite. In an eigenbasis of $A$, (F9) is
the entrywise product of $U^*HU$ with $L_p(\lambda)$. Since the values and
derivatives of $p,f$ agree at every eigenvalue, $L_p(\lambda)=L_f(\lambda)$.
Combining (F8)--(F9) proves
\[
 Df(A)[H]=U\bigl(L_f(\lambda)\circ(U^*HU)\bigr)U^*.
\]
In fact, differentiability of $f$ at the finitely many eigenvalues, plus
being defined near them, suffices for this pointwise theorem. The catalogue's
$C^1$ hypothesis is retained for its surrounding API. Complex matrix entries
do not turn this into a complex derivative of an arbitrary real function.

\paragraph{Verified reuse.}
The existing real Weyl perturbation theorem can supply (F8), but the unit
eigenvector proof above also covers complex Hermitian matrices and needs no
sorting or continuous choice of eigenvectors. The exact existing bridge
\code{NLAlib.specNorm_eq_norm} holds for all \code{RCLike} rectangular
matrices. The polynomial divided-difference lemmas from the preceding
subsection identify both the confluent and nonconfluent entries of (F9).

\subsection{L\"owner's criterion with a scalar mean-value proof}
\label{sec:matrix-monotone-loewner-criterion}
\textbf{Atlas:} \code{matrix-monotone-loewner-criterion}.
For $f\in C^1(I)$, matrix monotonicity in every size is equivalent to
$L_f(\lambda)\succeq0$ for every finite list $\lambda\in I^n$.

Necessity: at $A=\diag(\lambda)$ use the PSD direction $J=\boldsymbol1
\boldsymbol1^*$. For all sufficiently small $t>0$, the spectrum of $A+tJ$
remains in $I$ and
$[f(A+tJ)-f(A)]/t\succeq0$. Its limit is
$Df(A)[J]=L_f(\lambda)$. The PSD cone is closed, as is seen by taking the
limit of each quadratic form. Repeated nodes require no perturbation.

Sufficiency: let $A\preceq B$ and put $D=B-A\succeq0$.
Rayleigh quotients keep $A+tD$ in a compact spectral subinterval of $I$
for $0\le t\le1$. In any eigenbasis, its derivative in direction $D$
is a Schur product of two PSD matrices, hence is PSD. A finite Gram proof
of this fact is $K_{ij}=\langle u_i,u_j\rangle$,
$D_{ij}=\langle v_i,v_j\rangle$, so
$(K\circ D)_{ij}=\langle u_i\otimes v_i,u_j\otimes v_j\rangle$.
For every fixed vector $v$, the real function
$t\mapsto v^*f(A+tD)v$ is differentiable and has nonnegative derivative.
The ordinary real mean-value theorem makes it nondecreasing. Evaluating at
$0,1$ proves $f(A)\preceq f(B)$. This avoids integrating a matrix-valued
derivative and does not need a continuous eigenbasis along the segment.
The quantifier over all finite dimensions is part of the assertion.

\subsection{The noncommutative exponential derivative with an explicit remainder}
\label{sec:matrix-exponential-frechet-integral}
\textbf{Atlas:} \code{matrix-exponential-frechet-integral}.
For arbitrary real or complex square matrices,
\[
 D\exp(A)[E]=\int_0^1e^{(1-t)A}Ee^{tA}\,dt.                    \tag{F10}
\]
The curve $F(t)=e^{(1-t)(A+E)}e^{tA}$ has derivative
$F'(t)=-e^{(1-t)(A+E)}Ee^{tA}$. Only differentiation of a scalar multiple
of a fixed matrix is used: each matrix commutes with its own exponential.
The fundamental theorem gives the exact identity
\[
 e^{A+E}-e^A=\int_0^1e^{(1-t)(A+E)}Ee^{tA}\,dt.
\]
The telescoping formula for powers and the exponential series give
\[
 \norm{e^U-e^V}\le\norm{U-V}e^{\max(\norm U,\norm V)}.
\]
Subtracting (F10) from the exact identity therefore gives
\[
 \left\|e^{A+E}-e^A-\int_0^1e^{(1-t)A}Ee^{tA}\,dt\right\|
 \le\tfrac12 e^{\norm A+\norm E}\norm E^2.                    \tag{F11}
\]
The proposed derivative is bounded and linear, with norm at most
$e^{\norm A}$. This proves Fr\'echet differentiability on the full algebra,
with no normality or commutativity condition.

\paragraph{Verified reuse.}
Mathlib's \code{hasDerivAt_exp_smul_const} has conclusion
$\mathrm{HasDerivAt}(t\mapsto\exp(t\,x))\,(\exp(t\,x)x)\,t$
in a complete normed algebra, without assuming that algebra commutative.
Its primed variant places $x$ on the left. These are the appropriate tools
for the curve $F$. The unrestricted-looking
\code{hasFDerivAt_exp} is in the commutative-algebra section
and is not the theorem to use on arbitrary matrices.
\code{NLAlib.matrixExp_smul_eq_cfc} supplies the existing Hermitian bridge
$\mathrm{matrixExp}(sA)=\mathrm{cfc}(x\mapsto e^{sx})(A)$.

\subsection{Geometric mean and its maximal block property}
\label{sec:matrix-geometric-mean-def}
\textbf{Atlas:} \code{matrix-geometric-mean-def},
\code{matrix-geometric-mean-block-maximum}.
For $A\succ0$, $B\succeq0$, set
\[
 C=A^{-1/2}BA^{-1/2},\qquad G=A\#B=A^{1/2}C^{1/2}A^{1/2}.
\]
Then $G\succeq0$, $GA^{-1}G=B$, and $A\#A=A$ by multiplication.
For any Hermitian $X$, congruence by
\label{sec:matrix-geometric-mean-block-maximum}
$\diag(A^{-1/2},A^{-1/2})$ and completion of the square show
\[
 \begin{bmatrix}A&X\\X&B\end{bmatrix}\succeq0
 \quad\Longleftrightarrow\quad C-Y^2\succeq0,
 \qquad Y=A^{-1/2}XA^{-1/2}.
\]
Indeed the transformed quadratic form is
$\norm{u+Yv}^2+v^*(C-Y^2)v$. The spectral theorem gives
$Y\preceq |Y|=(Y^2)^{1/2}$, and square-root monotonicity gives
$|Y|\preceq C^{1/2}$. Congruence back yields $X\preceq G$.
Taking $Y=C^{1/2}$ attains equality and proves block positivity at $G$.
There is no positivity assumption on $X$.

\paragraph{Verified reuse.}
\code{CFC.sqrt_le_sqrt} in Mathlib's
\code{Analysis.SpecialFunctions.ContinuousFunctionalCalculus.Rpow.Order}
states $a\le b\Rightarrow\sqrt a\le\sqrt b$ for a nonunital
$C^*$-algebra with its star order. It covers complex matrices, so the
real-only theorem \code{NLAlib.loewner_heinz_real} need not be silently
extended. The real statement can also follow by complexifying a real
symmetric matrix and observing that its square root is real.
The domain $A\succ0$ is essential to this formula. A singular extension
is not obtained by replacing the inverse by a pseudoinverse.

\subsection{Even Schatten smoothness from a diagonal Hessian}
\label{sec:schatten-even-moment-smoothness}
\textbf{Atlas:} \code{schatten-even-moment-smoothness}.
Let $p=2r\ge2$. For Hermitian $A,H$, put $F(A)=\tr(A^p)\ge0$.
Differentiating the finite polynomial twice, and diagonalizing only $A$,
gives
\[
 D^2F(A)[H,H]=p\sum_{i,j}|h_{ij}|^2
       \sum_{q=0}^{p-2}\lambda_i^q\lambda_j^{p-2-q}.             \tag{F12}
\]
For $p>2$, weighted scalar AM--GM, followed by summation in $q$, gives
\[
 \sum_{q=0}^{p-2}\lambda_i^q\lambda_j^{p-2-q}
 \le\frac{p-1}{2}(|\lambda_i|^{p-2}+|\lambda_j|^{p-2}).
\]
Thus
\[
 D^2F(A)[H,H]\le p(p-1)\sum_i|\lambda_i|^{p-2}(H^2)_{ii}
 \le p(p-1)F(A)^{1-2/p}\norm H_{S_p}^2.                       \tag{F13}
\]
The last inequality uses only finite scalar H\"older and
\[
 \sum_i((H^2)_{ii})^{p/2}\le\tr((H^2)^{p/2})=\norm H_{S_p}^{p}.
\]
For completeness, diagonalize $H^2=V\diag(\eta_j)V^*$. Each diagonal
entry is the convex combination $\sum_j|V_{ij}|^2\eta_j$; scalar Jensen,
then the column sums of $|V_{ij}|^2$, proves the displayed inequality.
For $p=2$, (F13) is directly $D^2F(A)[H,H]=2\norm H_{S_2}^2$.

For $\varepsilon>0$, along any affine Hermitian line set
$g_\varepsilon(t)=(F(A+tH)+\varepsilon)^{2/p}$. The second derivative
has one term with coefficient $(2/p)(2/p-1)\le0$. Discarding that term
and using (F13) proves
\[
                     g_\varepsilon''(t)\le2(p-1)\norm H_{S_p}^2.
                                                                    \tag{F14}
\]
Regularization avoids every division by $F=0$.

Now apply (F14) to the block-diagonal line
\[
 Z(t)=\diag(X+tY,X-tY),\qquad H=\diag(Y,-Y),
\]
initially for Hermitian $X,Y$. The function is even, so
$g_\varepsilon'(0)=0$. Two applications of the real fundamental theorem
give $g_\varepsilon(1)\le g_\varepsilon(0)+(p-1)\norm H_{S_p}^2$.
Let $\varepsilon\downarrow0$ and cancel $2^{2/p}$ to obtain
\[
 \left(\frac{\norm{X+Y}_{S_p}^p+\norm{X-Y}_{S_p}^p}{2}\right)^{2/p}
 \le\norm X_{S_p}^2+(p-1)\norm Y_{S_p}^2.                      \tag{F15}
\]
For rectangular real or complex $X,Y$, apply this to the Hermitian
dilations $\left[\begin{smallmatrix}0&X\\X^*&0\end{smallmatrix}\right]$.
Their $p$th trace moments equal twice those of the original singular values;
the same factor cancels on both sides. Zero-sized rectangles are immediate.

\paragraph{Why this is a smaller formalization.}
Only polynomial differentiation, the spectral theorem, scalar finite
H\"older/Jensen, and one-variable calculus enter. There is no generalized
Klein input, no change of variable $t\mapsto\sqrt t$, and no complex
extension of a real-only Schatten H\"older theorem. The verified
\code{NLAlib.generalized_klein} remains an alternative: it already lifts
finite scalar inequalities to trace inequalities over \code{RCLike}.
Both \code{NLAlib.schattenNorm} and \code{NLAlib.schattenNormReal} are
defined over \code{RCLike}, although several currently indexed norm
inequalities still have real-only signatures.

\subsection{Rational block differences by finite triangular inversion}
\label{sec:matrix-function-block-divided-differences}
\textbf{Atlas:} \code{matrix-function-block-divided-differences}.
Work in the full complex matrix algebra. If $f=P/Q$ and all relevant
$Q(A)$ are invertible, define $f(A)=P(A)Q(A)^{-1}$. Define $\Delta f(A,B)[H]$
and $\Delta^2 f(A,B,C)[G,H]$ as the $(1,2)$ and $(1,3)$ blocks of $f$ on
the respective triangular matrices
\[
 \begin{bmatrix}A&H\\0&B\end{bmatrix},\qquad
 \begin{bmatrix}A&G&0\\0&B&H\\0&0&C\end{bmatrix}.
\]
For the monomial $z^m$, multiplication of block matrices gives the finite
formulas
\begin{align*}
 \Delta z^m(A,B)[H]&=\sum_{a+b=m-1}A^aHB^b,\\
 \Delta^2 z^m(A,B,C)[G,H]&=\sum_{a+b+c=m-2}A^aGB^bHC^c.
\end{align*}
Empty sums are zero. The first expression is linear and the second bilinear.
For general products the same block multiplication gives
\begin{align*}
 \Delta(fg)&=f(A)\Delta g(A,B)+\Delta f(A,B)g(B),\\
 \Delta^2(fg)&=f(A)\Delta^2g(A,B,C)
 +\Delta f(A,B)\Delta g(B,C)+\Delta^2f(A,B,C)g(C),              \tag{F16}
\end{align*}
where the middle term takes directions $G,H$ in that order. For $q^{-1}$,
write $q_D=q(D)$; triangular inversion gives
\begin{align*}
 \Delta(q^{-1})(A,B)[H]&=-q_A^{-1}\Delta q(A,B)[H]q_B^{-1},\\
 \Delta^2(q^{-1})(A,B,C)[G,H]
 &=q_A^{-1}\Delta q(A,B)[G]q_B^{-1}\Delta q(B,C)[H]q_C^{-1}\\
 &\quad-q_A^{-1}\Delta^2q(A,B,C)[G,H]q_C^{-1}.                 \tag{F17}
\end{align*}
In particular, invertibility of the diagonal blocks is sufficient; no
extra spectral hypothesis on a nonnormal triangular matrix is needed.

These formulas give, by telescoping monomials and then the product/inverse
rules (F16)--(F17),
\begin{align*}
 f(A)-f(B)&=\Delta f(A,B)[A-B],\\
 \Delta f(A,B)[H]-\Delta f(A,C)[H]
 &=\Delta^2 f(A,B,C)[H,B-C],\\
 \Delta f(A,C)[H]-\Delta f(B,C)[H]
 &=\Delta^2 f(A,B,C)[A-B,H].                                \tag{F18}
\end{align*}
A useful direct way to verify the last two identities, avoiding a large
rational induction, is a block intertwiner. For the first of them let
$M=\left[\begin{smallmatrix}A&H&0\\0&B&B-C\\0&0&C\end{smallmatrix}\right]$.
Its invariant embeddings of $(u,v)$ as $(u,v,0)$ and $(u,v,-v)$ have
restrictions $\left[\begin{smallmatrix}A&H\\0&B\end{smallmatrix}\right]$
and $\left[\begin{smallmatrix}A&H\\0&C\end{smallmatrix}\right]$.
Applying $f(M)$ and taking the first component gives that identity. For
the second, use
$M=\left[\begin{smallmatrix}A&A-B&0\\0&B&H\\0&0&C\end{smallmatrix}\right]$;
the quotient maps $(x,y,z)\mapsto(y,z)$ and
$(x,y,z)\mapsto(x+y,z)$ intertwine $M$ respectively with the $(B,C)$
and $(A,C)$ two-block matrices. Taking their off-diagonal blocks gives
the claimed difference. Polynomial intertwiners extend to rational ones
by multiplying by the inverses of the respective $Q$ evaluations.

The domain $\{A:Q(A)\text{ invertible}\}$ is open. Polynomial
differentiation and $D(V^{-1})[H]=-V^{-1}HV^{-1}$ give
\[
 Df(A)[H]=\Delta f(A,A)[H],\qquad
 D^2f(A)[G,H]=\Delta^2f(A,A,A)[G,H]+\Delta^2f(A,A,A)[H,G].      \tag{F19}
\]
For example, the inversion derivative follows from the exact resolvent
identity, with the remaining term quadratic in the perturbation; repeating
the product rule gives both orders in (F19). Thus full Fr\'echet derivatives,
not merely directional derivatives, are obtained. With
$R_\zeta(A)=(\zeta I-A)^{-1}$, (F17) gives
\[
 \Delta^2R_\zeta(A,B,C)[G,H]=R_\zeta(A)G R_\zeta(B)H R_\zeta(C).
\]
The ordered second difference is generally not half the Hessian: for
$f(z)=z^2$ it is $GH$, while half the Hessian is $(GH+HG)/2$.
The half-Hessian statement is valid for equal directions. None of these
triangular evaluations uses Hermitian continuous functional calculus.

\subsection{K-transform stability without a Perron--Frobenius theorem}
\label{sec:matrix-k-transform-stability}
\textbf{Atlas:} \code{matrix-k-transform-stability}.
Let $k\ge1$, $\alpha>0$, and $\nu$ be a compactly supported probability
measure on Hermitian $k\times k$ matrices. Assume $S\succ0$ and
$S^{-1}+T\succ0$ for every $T$ in its support. Set
\[
 K(S)=\int T(I+ST)^{-1}\,d\nu(T)-\alpha^{-1}S^{-1},\quad
 Z_T=S^{1/2}TS^{1/2},\quad C_T=Z_T(I+Z_T)^{-1},
\]
and $\mathcal L(H)=\alpha\int C_THC_T\,d\nu(T)$ on the real inner-product
space of Hermitian matrices, with $\langle H,J\rangle_F=\tr(HJ)$.
The matrices $C_T$ are Hermitian, since they are real rational functions of
$Z_T$ and $I+Z_T\succ0$. Hence $\mathcal L$ preserves PSD matrices and is
self-adjoint for the Frobenius inner product.

\paragraph{Derivative and congruence.}
Inversion differentiation gives
\[
 DK(S)[\Delta]=\alpha^{-1}S^{-1}\Delta S^{-1}
 -\int T(I+ST)^{-1}\Delta(I+TS)^{-1}T\,d\nu(T).               \tag{F20}
\]
Here $T(I+ST)^{-1}=(I+TS)^{-1}T
=S^{-1/2}C_TS^{-1/2}$, so, with $H=S^{-1/2}\Delta S^{-1/2}$,
\[
 DK(S)[\Delta]=\alpha^{-1}S^{-1/2}(H-\mathcal L(H))S^{-1/2},
 \quad
 \langle\Delta,DK(S)[\Delta]\rangle_F
 =\alpha^{-1}\langle H,(I-\mathcal L)H\rangle_F.              \tag{F21}
\]
The derivative under the integral is legitimate with the stated hypotheses:
compactness supplies a positive minimum of
$\lambda_{\min}(S^{-1}+T)$ and a finite maximum of $\norm T$ on the
support. On a sufficiently small neighbourhood of $S$, continuity of
inversion retains half this gap uniformly. The inverse identity gives a
quadratic remainder bounded by one constant times $\norm\Delta^2$
uniformly in $T$. Integrating this finite bound proves (F20) as a
Fr\'echet derivative. No global assumption on $T$ outside the support is
required.

\paragraph{A finite spectral lemma for cone-positive maps.}
If $L$ is self-adjoint on Hermitian matrices and preserves the PSD cone,
then
\[
                     \norm L_{F\to F}=\lambda_{\max}(L)\ge0. \tag{F22}
\]
Indeed write $H=H_+-H_-$, $|H|=H_++H_-$. Every pairing
$\langle H_\pm,LH_\pm\rangle_F$ and
$\langle H_+,LH_-\rangle_F$ is nonnegative, because the trace of a
product of two PSD matrices is nonnegative. Self-adjointness therefore gives
\[
 |\langle H,LH\rangle_F|\le\langle|H|,L|H|\rangle_F,
 \qquad\norm{|H|}_F=\norm H_F.
\]
Rayleigh maximization for the finite real symmetric operator proves (F22).
This does \emph{not} claim that $L$ is a PSD operator for the Frobenius
inner product.

Combining (F21)--(F22), $DK(S)$ is Frobenius positive definite if and only
if $\norm{\mathcal L}_{F\to F}<1$. The latter implies convergence of the
Neumann series; each term preserves the PSD cone, so
$(I-\mathcal L)^{-1}=\sum_{j\ge0}\mathcal L^j$ does too. In particular
\[
 Y=(I-\mathcal L)^{-1}I\succeq I\succ0,
 \qquad\mathcal L(Y)=Y-I\prec Y.                            \tag{F23}
\]

Conversely, if $Y\succ0$ and $\mathcal L(Y)\prec Y$, then for some
$0\le q<1$ one has $0\preceq\mathcal L(Y)\preceq qY$.
For every Hermitian $H$, choose $c$ with $-cY\preceq H\preceq cY$.
Positivity and induction give
$-cq^jY\preceq\mathcal L^j(H)\preceq cq^jY$.
Apply this to a nonzero eigenvector of the self-adjoint map $\mathcal L$:
norm/order equivalence gives $|\lambda|^j\norm H\le cq^j\norm Y$,
so $|\lambda|\le q$. If $q=0$, the $j=1$ inequality directly makes
$\mathcal L=0$. Thus its Frobenius norm is less than one.
Finally (F21) gives the equivalence of (F23)'s strict supersolution with
\[
 \exists B\succ0:\ DK(S)[B]\succ0,
 \qquad B=S^{1/2}YS^{1/2}.
\]
This proves all four stated equivalences.

\paragraph{Correction retained.}
Cone positivity must not be confused with a nonnegative Hilbert quadratic
form. For example take $S=I$, $\alpha=1$, and $\nu$ concentrated at
$T=\diag(1,-1/2)$. Then $I+T\succ0$ and
$C=\diag(1/2,-1)$. For
$H=\left[\begin{smallmatrix}0&1\\1&0\end{smallmatrix}\right]$,
$CHC=-H/2$, so $\langle H,\mathcal L(H)\rangle_F<0$.
The replacement (F22) is precisely what repairs the argument.

\subsection{Coefficient block partial transpose as a swap compression}
\label{sec:coefficient-block-partial-transpose}
\textbf{Atlas:} \code{coefficient-block-partial-transpose}.
For Hermitian $A_1,\ldots,A_s\in\C^{d\times d}$ with $s,d\ge1$, define
\[
 C_{ij}=A_jA_i,\quad \sigma^2=\left\|\sum_iA_i^2\right\|,
 \quad v^2=\left\|\sum_i\operatorname{vec}(A_i)
                 \operatorname{vec}(A_i)^*\right\|.
\]
Since $C_{ji}^*=C_{ij}$, $C$ is Hermitian. Its diagonal blocks are $A_i^2$,
so any unit vector supported in one block has a nonnegative Rayleigh
quotient. Thus $\lambda_{\max}(C)\ge0$ and
$\kappa=\sqrt{\lambda_{\max}(C)}$ is meaningful.

For $x=(x_1,\ldots,x_s)$, finite Cauchy--Schwarz gives
\begin{align*}
 |x^*Cx|&=\left|\sum_{i,j}\langle A_jx_i,A_ix_j\rangle\right|\\
 &\le\sum_{i,j}\norm{A_jx_i}^2
 =\sum_i x_i^*\left(\sum_jA_j^2\right)x_i
 \le\sigma^2\norm x^2.
\end{align*}
The two squared-norm sums in Cauchy--Schwarz are equal after interchanging
$i,j$. Hermitian norm/Rayleigh equivalence proves $\norm C\le\sigma^2$.

For the covariance bound define the finite matrix
$T:\C^s\to\C^d\otimes\C^d$ by
\[
 Te_i=\sum_{a,b}(A_i)_{ab}e_a\otimes e_b,
 \qquad F=T\otimes I_d.
\]
Let $S_{13}(u\otimes v\otimes w)=w\otimes v\otimes u$.
It is an orthogonal coordinate permutation and $\norm{S_{13}}=1$.
Directly, with rows and columns indexed by $(i,c),(j,e)$,
\[
 (F^*S_{13}F)_{(i,c),(j,e)}
 =\sum_b\overline{(A_i)_{eb}}(A_j)_{cb}
 =\sum_b(A_j)_{cb}(A_i)_{be}=(A_jA_i)_{ce}.
\]
Therefore
\[
                 C=F^*S_{13}F,\qquad
 \norm C\le\norm F^2=\norm T^2=\norm{TT^*}=v^2.             \tag{F24}
\]
The equality $\norm{T\otimes I_d}=\norm T$ follows by reindexing into
$d$ copies of $T$, or by singular values; no abstract infinite tensor
product is needed. A different fixed vectorization convention only
conjugates $TT^*$ by a coordinate permutation.

Combining the bounds gives
$\norm C\le\min(\sigma^2,v^2)$ and
$\kappa\le\min(\sigma,v)$, where $\sigma,v$ denote the nonnegative
square roots of the displayed squared parameters. If the $A_i$ commute,
simultaneous unitary diagonalization writes their eigenvalues as real
numbers $a_{i\ell}$. Reindexing the block space gives
\[
 C\cong\bigoplus_{\ell=1}^d a_\ell a_\ell^*,
 \qquad a_\ell=(a_{1\ell},\ldots,a_{s\ell})^\top.
\]
Hence $\lambda_{\max}(C)=\max_\ell\sum_i a_{i\ell}^2=\sigma^2$,
so $\kappa=\sigma$. The finite simultaneous diagonalization itself follows
by restricting all other commuting Hermitian matrices to the eigenspaces
of the first and inducting on the finite family.

Finally, replacing $A_i$ by $A_i\otimes I_r$, $r\ge1$, produces $C\otimes I_r$
up to a coordinate permutation. Every eigenvalue of $C$ is repeated $r$
times, so its largest eigenvalue and $\kappa$ are unchanged. The block
$C$ need not be PSD; (F24) is a compression of a swap, which has both
signs. This is why $\kappa$ uses $\lambda_{\max}(C)$ rather than
$\norm C$.

\paragraph{Sources and scope.}
The catalogue contracts for the positive-map, derivative, mean, and
smoothness entries cite
\href{https://tropp.caltech.edu/notes/Tro22-Matrix-Analysis-LN.pdf}{Tropp's
Matrix Analysis notes}. The two latest contracts come from
\href{https://arxiv.org/abs/2606.27774v1}{Montanari--Saeed, version 1}
and \href{https://arxiv.org/abs/2610.01982v1}{Bandeira et al., version 1}.
The proofs above explicitly supply the finite reductions and retain the
stated hypotheses; no Fock-space formalization or matrix-law asymptotics
are prerequisites for these fifteen entries.


% Matrix equations, regularization, and elementary matrix functions.
% Audited against the site snapshots dated 2026-10-15.
\section{Regularization and matrix equations from finite coordinates}
\label{sec:eq-frontier}

Throughout this section $\mathbb K$ is $\R$ or $\C$, and $*$ denotes the
adjoint. All vector norms are Euclidean. Empty orthonormal families and
$0\times0$ identity matrices are permitted; assertions involving a smallest
or largest eigenvalue explicitly require a nonzero dimension. Singular
values are numbered starting at one in the mathematics. A fixed SVD is
used at repeated singular values, so no canonical choice of a truncated
subspace is asserted.

\paragraph{Reusable coordinate lemma.}
Let $u_1,\ldots,u_r$ and $v_1,\ldots,v_r$ be orthonormal families and let
$A=\sum_{i=1}^r\sigma_i u_i v_i^*$, with $\sigma_i>0$. Then
\[
 A^*A v_i=\sigma_i^2v_i,\qquad
 A^\dagger=\sum_{i=1}^r\sigma_i^{-1}v_i u_i^*,\qquad
 \ker A=\operatorname{span}\{v_i\}_{i=1}^r{}^\perp.
\]
For any scalar coefficients $a_i$,
\[
 \left\|\sum_i a_i u_i v_i^*\right\|_2=\max_i|a_i|
 \quad(r>0),\qquad
 \left\|\sum_i a_i u_i v_i^*\right\|_F^2=\sum_i|a_i|^2.
\]
The first equality follows by applying the matrix to
$x=\sum_i\xi_i v_i+x_\perp$: the squared output norm is
$\sum_i|a_i|^2|\xi_i|^2$, and equality is attained at a maximizing $v_i$.
For $r=0$ both matrices and norms are zero. These identities also prove
the displayed pseudoinverse formula directly from the four Penrose
identities. They are the only SVD calculations needed below.

\paragraph{Verified reuse.}
The current declaration index supplies
\code{NLAlib.exists_isSVD},
\code{NLAlib.IsSVD.isMoorePenroseInverse},
\code{NLAlib.moorePenroseInverse}, and
\code{NLAlib.specNorm_eq_norm}, all over \code{RCLike} and with
\code{sorry=false}. The indexed least-squares normal-equation theorem
\code{NLAlib.isLeastSquaresSolution_iff_transpose_mulVec_mulVec_eq}
is real-valued. The completed-square identities below provide the small
complex extension directly, rather than assuming that theorem already has
the complex interface.

\subsection{Tikhonov filtering by one completed square}
\label{sec:eq-tikhonov}
\textbf{Atlas id:} \code{tikhonov-spectral-filter}.

Let $A:\mathbb K^n\to\mathbb K^m$, $b\in\mathbb K^m$, and $\beta>0$.
Put $H=A^*A+\beta I$. For every nonzero $z$,
\[
 z^*Hz=\|Az\|^2+\beta\|z\|^2>0,
\]
so $H$ is invertible. Set $x_\beta=H^{-1}A^*b$. Expanding at
$x_\beta+z$ and using $Hx_\beta=A^*b$ gives the exact identity
\begin{equation}\label{eq:eq-tikh-square}
 \|A(x_\beta+z)-b\|^2+\beta\|x_\beta+z\|^2
 =\|Ax_\beta-b\|^2+\beta\|x_\beta\|^2+z^*Hz.
\end{equation}
Thus $x_\beta$ is the unique minimizer, without any rank assumption.
The shared coordinate lemma gives
\[
 x_\beta=\sum_{i=1}^r
 \frac{\sigma_i}{\sigma_i^2+\beta}\,v_i u_i^*b.
\]
There is no kernel term because $A^*b\perp\ker A$. Every displayed
coefficient tends to $\sigma_i^{-1}$ as $\beta\downarrow0$, and a finite
sum therefore yields $x_\beta\to A^\dagger b$. Finally,
$(\sigma-\sqrt\beta)^2\ge0$ gives
$\sigma/(\sigma^2+\beta)\le1/(2\sqrt\beta)$, whence
\[
 \|x_\beta(b+e)-x_\beta(b)\|
 \le\frac{\|e\|}{2\sqrt\beta}.
\]
This proof works also for $A=0$ and empty domain or codomain: the relevant
finite sum is empty. No limiting inverse at $\beta=0$ is used.

\subsection{Truncated SVD: an exact orthogonal bias--noise identity}
\label{sec:eq-tsvd}
\textbf{Atlas id:} \code{truncated-svd-regularization}.

Fix $1\le k\le r=\rank A$ and write
$V_k=[v_1,\ldots,v_k]$, $U_k=[u_1,\ldots,u_k]$,
$\Sigma_k=\diag(\sigma_1,\ldots,\sigma_k)$. Every vector in
$\operatorname{range}V_k$ has the unique form $V_kz$. Since
$AV_k=U_k\Sigma_k$, orthogonal decomposition of $b$ gives
\[
 \|AV_kz-b\|^2
 =\|\Sigma_kz-U_k^*b\|^2+\|(I-U_kU_k^*)b\|^2.
\]
All $\sigma_i$ in this block are positive, so the unique minimizer is
$z=\Sigma_k^{-1}U_k^*b$. Consequently
\[
 x_k(b)=V_k\Sigma_k^{-1}U_k^*b=A_k^\dagger b,
 \qquad \|A_k^\dagger\|_2=\sigma_k^{-1}.
\]
For $b=Ax+e$, put $P=V_kV_k^*$. The exact error decomposition is
\[
 x_k(b)-x=-(I-P)x+A_k^\dagger e.
\]
Its two terms lie in orthogonal subspaces, giving the stronger identity
\[
 \|x_k(b)-x\|^2=\|(I-P)x\|^2+\|A_k^\dagger e\|^2
 \le\|(I-P)x\|^2+\sigma_k^{-2}\|e\|^2.
\]
For the optional $k=0$ convention, $x_0=0$ and the restricted domain is
$\{0\}$; the norm formula using $\sigma_k^{-1}$ must then be omitted.

\subsection{Landweber: solve a finite family of scalar recurrences}
\label{sec:eq-landweber}
\textbf{Atlas id:} \code{landweber-spectral-filter}.

Let $A\ne0$, $0<\eta<2/\|A\|_2^2$, and
$x_{t+1}=x_t+\eta A^*(b-Ax_t)$. Put
$a_i=1-\eta\sigma_i^2$ and $b_i=u_i^*b$. Taking coordinates gives
\[
 v_i^*x_{t+1}=a_i v_i^*x_t+\eta\sigma_i b_i,
 \qquad P_{\ker A}x_{t+1}=P_{\ker A}x_t.
\]
An induction, or the finite geometric-sum identity, gives
\begin{equation}\label{eq:eq-landweber}
 x_t=P_{\ker A}x_0+
 \sum_{i=1}^r v_i\left(
 a_i^t v_i^*x_0+\frac{1-a_i^t}{\sigma_i}b_i\right).
\end{equation}
The step restriction implies $-1<a_i<1$ for every positive singular
value. Hence $a_i^t\to0$ and
$x_t\to A^\dagger b+P_{\ker A}x_0$. Setting $x_0=0$ gives the card's
filter formula. More explicitly, with
$q=\max_{1\le i\le r}|a_i|<1$,
\[
 \|x_t-A^\dagger b-P_{\ker A}x_0\|
 \le q^t\|(I-P_{\ker A})x_0-A^\dagger b\|.
\]
This needs no consistency assumption on $b$. If $A=0$, the recurrence
itself gives $x_t=x_0$ for every $\eta$; the step restriction is not
interpreted by division by zero.

\subsection{Rectangular polar Newton iteration}
\label{sec:eq-polar-newton}
\textbf{Atlas id:} \code{polar-newton-iteration}.

Suppose $m\ge n\ge1$ and $A$ has full column rank. Choose a thin SVD
$A=U\diag(\sigma_i)V^*$, where $U^*U=I$, $V$ is unitary, and every
$\sigma_i>0$. Its polar isometry is $Q=UV^*$, since
$(A^*A)^{-1/2}=V\diag(\sigma_i^{-1})V^*$.
Inductively, the recurrence
\[
 X_0=A,\qquad X_{k+1}=\tfrac12\bigl(X_k+(X_k^\dagger)^*\bigr)
\]
has $X_k=U\diag(s_{i,k})V^*$ with
$s_{i,0}=\sigma_i$ and
$s_{i,k+1}=(s_{i,k}+s_{i,k}^{-1})/2>0$. Thus every iterate has full
column rank. The scalar change of variables
\[
 c(s)=\frac{s-1}{s+1},\qquad
 c\left(\frac{s+s^{-1}}2\right)=c(s)^2
\]
is verified after multiplying by the positive denominators. It follows
for all $k\ge0$ that
\[
 s_{i,k}=\frac{1+c_i^{2^k}}{1-c_i^{2^k}},\qquad
 c_i=\frac{\sigma_i-1}{\sigma_i+1},\qquad |c_i|<1.
\]
With $r=\max_i|c_i|<1$, the coordinate lemma gives
\[
 \|X_k-Q\|_2=\max_i|s_{i,k}-1|
 \le\frac{2r^{2^k}}{1-r^{2^k}}\longrightarrow0.
\]
The index $i$ denotes a fixed singular direction; no preservation of
decreasing order is needed. For $n=0$ the iteration is the unique empty
matrix and convergence is immediate, but the maximum defining $r$ is
omitted. Full rank is essential for the asserted isometry target.

\subsection{Inverse-free polar iteration: a scalar interval invariant}
\label{sec:eq-polar-schulz}
\textbf{Atlas id:} \code{polar-newton-schulz-iteration}.

Keep the full-column-rank assumption and assume additionally
$\|A\|_2<\sqrt3$. For
$X_{k+1}=\tfrac12X_k(3I-X_k^*X_k)$ the same thin SVD is preserved, with
\[
 s_{i,k+1}=f(s_{i,k}),\qquad f(s)=\tfrac12s(3-s^2).
\]
For $0<s<\sqrt3$ we have $f(s)>0$ and
\[
 1-f(s)=\tfrac12(s-1)^2(s+2)\ge0.
\]
Thus $0<s_{i,1}\le1$, and the interval $(0,1]$ is invariant.
Set $E_k=I-X_k^*X_k$. Direct polynomial expansion, using only that
$X_k^*X_k$ commutes with its polynomials, gives
\begin{equation}\label{eq:eq-polar-defect}
 E_{k+1}=\tfrac14 E_k^2(3I+E_k).
\end{equation}
For $k\ge1$, its eigenvalues $e_{i,k}=1-s_{i,k}^2$ belong to $[0,1)$.
The scalar recurrence
$e_{i,k+1}=e_{i,k}^2(3+e_{i,k})/4\le e_{i,k}^2$ therefore yields
\[
 0\preceq E_k\prec I,\qquad
 \delta_{k+1}\le\delta_k^2,\qquad
 \delta_k\le\delta_1^{2^{k-1}}\longrightarrow0,
 \quad\delta_k=\|E_k\|_2.
\]
Also $1-s\le1-s^2$ on $(0,1]$, and hence
$\|X_k-Q\|_2\le\delta_k$ for $k\ge1$.
The upper endpoint must be strict: $f(\sqrt3)=0$. The initial defect
need not have norm below one, so the quadratic estimate begins at
$k=1$, exactly as stated in the card.

\subsection{A normalized GSVD by whitening only the active subspace}
\label{sec:eq-gsvd}
\textbf{Atlas id:} \code{gsvd-normalized-pair}.

Let $A:\mathbb K^n\to\mathbb K^m$ and
$L:\mathbb K^n\to\mathbb K^p$. Put
\[
 G=A^*A+L^*L,\qquad N=\ker A\cap\ker L.
\]
The identity $x^*Gx=\|Ax\|^2+\|Lx\|^2$ proves
$\ker G=N$. Consequently $\rank G=\rank[A;L]=q$.
Choose orthonormal basis matrices $V$ for $N^\perp$ and $W$ for $N$.
Then $[V\ W]$ is unitary and $G_0=V^*GV$ is positive definite if $q>0$.
Diagonalize $G_0$ and form its positive inverse square root $R=G_0^{-1/2}$.
The matrix
\[
 M=R V^*A^*AV R
\]
is Hermitian and satisfies $0\preceq M\preceq I$, because its quadratic
form is nonnegative and
$I-M=RV^*L^*LV R\succeq0$. Choose a unitary $T$ with
$T^*MT=\diag(\theta_1,\ldots,\theta_q)$, $0\le\theta_i\le1$.
Define
\[
 X=[VR T\ \ W],\qquad
 c_i=\sqrt{\theta_i},\quad s_i=\sqrt{1-\theta_i}\quad(i\le q),
 \qquad c_i=s_i=0\quad(i>q).
\]
The active block $RT$ is invertible, so $X$ is invertible. Direct
multiplication gives
\[
 X^*A^*AX=\diag(c_i^2),\qquad
 X^*L^*LX=\diag(s_i^2).
\]
All cross terms involving $W$ vanish because $AW=LW=0$.
Thus $c_i^2+s_i^2=1$ on the active coordinates, and the last columns
span exactly $N$. If $c_i>0$, put $u_i=Ax_i/c_i$. The Gram identity
gives $u_i^*u_j=\delta_{ij}$ whenever $c_i,c_j>0$; the corresponding
statement for $Lx_i/s_i$ is identical.

This construction includes $q=0$: then $A=L=0$, choose any orthonormal
basis $W$, and take $X=W$ with all $c_i=s_i=0$. Empty row spaces cause
no change. The ratios $c_i/s_i$ are used only when their meaning is
specified: $s_i=0<c_i$ is an infinite generalized singular value;
$c_i=s_i=0$ is a common-kernel coordinate, not a ratio.

\paragraph{Verified reuse.}
\code{NLAlib.exists_definite_pair_congruence_diagonal} already handles
Hermitian definite pairs over \code{RCLike}. It can replace the whitening
and diagonalization step on $N^\perp$. The more specialized
\code{NLAlib.exists_normalized_diagonal_congruence_of_posSemidef}
is currently real-valued; its normalization is different from
$c_i^2+s_i^2=1$. The construction above explicitly supplies this final
normalization and the common-kernel block.

\subsection{General-form Tikhonov and its two limits}
\label{sec:eq-general-tikhonov}
\textbf{Atlas id:} \code{general-form-tikhonov-gsvd}.

For $\beta>0$ let $H_\beta=A^*A+\beta L^*L$ and $h=A^*b$.
Its kernel is $N=\ker A\cap\ker L$. Moreover $h\perp N$, so
$h\in\operatorname{range}H_\beta$, by Hermitian diagonalization.
Thus $x_*=H_\beta^\dagger h$ solves $H_\beta x_*=h$.
Expanding the objective at $x_*+z$ gives
\[
 \|A(x_*+z)-b\|^2+\beta\|L(x_*+z)\|^2
 =\|Ax_*-b\|^2+\beta\|Lx_*\|^2+z^*H_\beta z.
\]
Its minimizers are exactly $x_*+N$, and uniqueness is equivalent to
$N=\{0\}$.

Assume now $N=\{0\}$ and use the normalized GSVD above, so $q=n$.
The congruence identity
$X^*H_\beta X=D_\beta=\diag(c_i^2+\beta s_i^2)$ gives
$H_\beta^{-1}=X D_\beta^{-1}X^*$, not a unitary-similarity formula.
Writing $x_\beta=Xy_\beta$, we obtain
\[
 (y_\beta)_i=\frac{g_i}{c_i^2+\beta s_i^2},\qquad
 g_i=x_i^*A^*b=
 \begin{cases}c_i u_i^*b,&c_i>0,\\0,&c_i=0.\end{cases}
\]
All denominators are positive.

As $\beta\downarrow0$, the coordinate limit is $u_i^*b/c_i$ when
$c_i>0$, and zero when $c_i=0$. To identify it variationally, expand
\[
 \|AXy-b\|^2
 =\sum_{c_i>0}|c_i y_i-u_i^*b|^2
   +\left\|b-\sum_{c_i>0}u_i u_i^*b\right\|^2,
 \qquad \|LXy\|^2=\sum_i s_i^2|y_i|^2.
\]
All least-squares minimizers have the displayed fixed coordinates for
$c_i>0$ and arbitrary coordinates for $c_i=0$. On the latter coordinates
$s_i=1$, so the unique smallest-penalty choice is zero. This proves the
small-parameter endpoint exactly.

As $\beta\to\infty$, the coordinates with $s_i>0$ tend to zero,
and those with $s_i=0$ are unchanged. Here $s_i=0$ implies $c_i=1$.
The condition $LXy=0$ is precisely $y_i=0$ for all $s_i>0$; minimizing
the remaining independent squares proves that this limit is the unique
least-squares minimizer constrained to $\ker L$. Both limits are finite
coordinate limits. In general the first is not $A^\dagger b$, since the
penalty is $L$, not the identity.

\subsection{Total least squares as a Rayleigh quotient on an open chart}
\label{sec:eq-tls}
\textbf{Atlas id:} \code{total-least-squares-svd}.

Let $m,n\ge1$, $C=[A\ b]$, and
$\mu=\lambda_{\min}(C^*C)$, with all $n+1$ domain-side eigenvalues
included. Write a data correction as $D=[E\ f]$. The constraint
$(A+E)x=b+f$ is equivalent to
\[
 (C+D)v=0,\qquad
 v=\frac{(x,-1)}{\sqrt{1+\|x\|^2}}.
\]
For any unit vector $v$, the exact minimum correction satisfying
$Dv=-Cv$ is
\begin{equation}\label{eq:eq-tls-correction}
 D_0=-Cvv^*,\qquad \|D_0\|_F^2=\|Cv\|^2.
\end{equation}
Indeed, every feasible $D$ has the orthogonal decomposition
$D=-Cvv^*+D(I-vv^*)$ in the Frobenius inner product. Thus its squared
norm is $\|Cv\|^2+\|D(I-vv^*)\|_F^2$, proving both the value and
uniqueness of the minimum for a fixed direction.

Vectors arising from finite $x$ are exactly the unit directions whose
last coordinate is nonzero, up to a nonzero scalar of modulus one.
This set is dense in the unit sphere: if the last coordinate of $v$
is zero, normalize $v+\varepsilon e_{n+1}$ and let $\varepsilon\downarrow0$.
The continuous Rayleigh quotient $v\mapsto v^*C^*Cv$ has sphere minimum
$\mu$, so its infimum over this dense chart is also $\mu$.
Equality $v^*C^*Cv=\mu$ is equivalent to
$v\in\ker(C^*C-\mu I)$, since diagonalization expresses the difference
as a sum of nonnegative squares. Consequently the TLS infimum is
attained if and only if this eigenspace has a vector with nonzero last
coordinate, and all optimal solutions are exactly
\[
 x=-v_{1:n}/v_{n+1}
 \quad\text{for such minimizing eigenvectors }v.
\]
For a unit choice, \eqref{eq:eq-tls-correction} gives the optimal data
correction. This includes repeated minimum eigenvalues; one must inspect
the eigenspace, not only a particular eigenbasis.

For the generic gap $\lambda_{\min}(A^*A)>\mu$, put
$P=A^*A-\mu I\succ0$. The first block of the eigenvector equation is
$Pv_{1:n}+A^*b\,v_{n+1}=0$. If $v_{n+1}=0$, it forces $v=0$.
Thus every nonzero vector in the minimizing eigenspace is feasible, and
the map taking its last coordinate is injective. The eigenspace is
therefore one-dimensional, and
\[
 x=P^{-1}A^*b
\]
is the unique optimum. If the gap fails, attainment must still be tested;
it cannot be inferred by totalized division by a zero last coordinate.
The argument uses only a Hermitian minimum and a rank-one correction,
not a general low-rank approximation theorem.

\subsection{Nearest PSD matrix: entrywise squares in one eigenbasis}
\label{sec:eq-nearest-psd}
\textbf{Atlas id:} \code{nearest-psd-frobenius}.

For arbitrary complex $A$, put $H=(A+A^*)/2$ and $K=(A-A^*)/2$.
Use the real inner product
$\langle Y,Z\rangle_{F,\R}=\operatorname{Re}\tr(Y^*Z)$.
Hermitian and skew-Hermitian matrices are orthogonal: the trace of their
product is purely imaginary, as follows by conjugating the trace and
cycling the two factors. Hence, for Hermitian $X$,
\[
 \|A-X\|_F^2=\|K\|_F^2+\|H-X\|_F^2.
\]
Choose $H=U\diag(\lambda_i)U^*$ and write $Y=U^*XU$. If
$X\succeq\delta I$, then $Y_{ii}\ge\delta$ and
\[
 \|H-X\|_F^2
 =\sum_i(\lambda_i-Y_{ii})^2+\sum_{i\ne j}|Y_{ij}|^2
 \ge\sum_i(\delta-\lambda_i)_+^2.
\]
Equality requires and is attained by
$Y=\diag(\max(\lambda_i,\delta))$. This proves uniqueness and the
exact distance formula for
$P_\delta(A)=U\diag(\max(\lambda_i,\delta))U^*$.

For completeness, nonexpansivity follows from an explicit variational
inequality. If $P=P_\delta(A)$ and $X=X^*\succeq\delta I$, then
\[
 \langle A-P,X-P\rangle_{F,\R}\le0.
\]
In the eigenbasis above, only the diagonal entries $\lambda_i<\delta$
contribute, and each contribution is
$(\lambda_i-\delta)(Y_{ii}-\delta)\le0$; $K$ contributes zero.
Apply this inequality once to $A$ with $X=P_\delta(B)$ and once to $B$
with $X=P_\delta(A)$. Adding gives
\[
 \|P_\delta(A)-P_\delta(B)\|_F^2
 \le\langle A-B,P_\delta(A)-P_\delta(B)\rangle_{F,\R}
 \le\|A-B\|_F\|P_\delta(A)-P_\delta(B)\|_F.
\]
If the last factor is zero the claim is immediate; otherwise divide by
it. This avoids a general metric-projection API. Empty dimensions are
handled by empty sums. The proof is valid for real matrices as well.

\paragraph{Verified reuse.}
The exact real-inner-product conventions and Cauchy--Schwarz bound are
provided by \code{NLAlib.frobInner_eq_re_trace_mul_conjTranspose} and
\code{NLAlib.frobInner_le_frobNorm_mul_frobNorm}, both over
\code{RCLike}. Unitary norm invariance follows from the same finite
sum identities already used in the SVD API.

\section{Matrix sign, residual squaring, and spectral backward error}
\label{sec:eq-functions}

\subsection{Matrix sign from two coprime factors, without Jordan form}
\label{sec:eq-sign-projectors}
\textbf{Atlas id:} \code{matrix-sign-spectral-projectors}.

Let $A\in\C^{n\times n}$ have no eigenvalue on the imaginary axis.
Factor its characteristic polynomial, retaining algebraic multiplicities,
as $p=p_+p_-$, with all roots of $p_+$ in the open right half-plane and
all roots of $p_-$ in the open left half-plane. A factor is $1$ when
its root set is empty. The factors are coprime, so choose polynomials
$a,b$ with $ap_++bp_-=1$. Define
\[
 P_+=b(A)p_-(A),\qquad P_-=a(A)p_+(A).
\]
By Cayley--Hamilton, $p_+(A)p_-(A)=0$. Therefore
$P_++P_-=I$, $P_+P_-=P_-P_+=0$, and $P_\pm^2=P_\pm$.
Moreover
\[
 \operatorname{range}P_+=\ker p_+(A),\qquad
 \operatorname{range}P_-=\ker p_-(A).
\]
For example, $p_+(A)P_+=0$, and if $p_+(A)v=0$ the Bezout identity
gives $P_+v=v$. Both subspaces are invariant under $A$.
Repeated application of the same coprime-factor argument splits
$\ker p_+(A)$ into the kernels of $(A-\lambda I)^{m_\lambda}$ for
right-half-plane roots, and similarly on the left. These are exactly
the corresponding generalized spectral subspaces; no Jordan basis is
required.

Set $S=P_+-P_-$. Then $S^2=I$ and $AS=SA$. On the two summands $S$
is respectively $I$ and $-I$. The restriction of $A$ to
$\ker p_+(A)$ has spectrum among the roots of $p_+$: an eigenvector
with eigenvalue $\lambda$ would satisfy $p_+(\lambda)v=0$.
The analogous left assertion shows that every eigenvalue of $SA$ has
positive real part.

To prove uniqueness, suppose $T^2=I$, $TA=AT$, and $TA$ is positive
stable. Since $T$ commutes with all polynomials in $A$, it preserves
both $\operatorname{range}P_\pm$. Each of these spaces splits into its
$T=1$ and $T=-1$ parts by the projectors $(I\pm T)/2$.
The intersection of $\operatorname{range}P_+$ with the $T=-1$ space
must be zero: otherwise the restriction of $A$ to this nonzero invariant
complex space has an eigenvalue $\lambda$ with $\operatorname{Re}\lambda>0$,
and $TA$ has eigenvalue $-\lambda$, a contradiction. Likewise the
intersection of $\operatorname{range}P_-$ with the $T=1$ space is zero.
Thus $T=S$. This proves the exact characterization and
$P_\pm=(I\pm S)/2$, including defective and nonnormal matrices.
These projectors need not be orthogonal. For $n=0$ the same polynomial
construction is vacuous and all identities hold.

\paragraph{Finite algebra to package.}
The only new reusable lemma is the coprime-kernel splitting
$\ker(f(A)g(A))=\ker f(A)\oplus\ker g(A)$ when $f,g$ are coprime.
Its explicit projections are the Bezout evaluations above. The
commutation and invertibility calculations then remain in the finite
commutative algebra generated by $A$.

\subsection{Newton sign iteration: define it first by the Cayley formula}
\label{sec:eq-sign-newton}
\textbf{Atlas id:} \code{matrix-sign-newton}.

Let $S=\operatorname{sign}(A)$ as just constructed. On a generalized
spectral subspace of $A$ with eigenvalue $\lambda$, $S=sI$ where
$s=1$ or $-1$ is the sign of $\operatorname{Re}\lambda$.
Thus $A+S$ is invertible. Define
$G=(A-S)(A+S)^{-1}$. Triangularizing $A$ separately on these invariant
subspaces shows that the eigenvalues of $G$ are
\[
 g_\lambda=\frac{\lambda-s}{\lambda+s}.
\]
The scalar identity
$|\lambda+s|^2-|\lambda-s|^2=4s\operatorname{Re}\lambda>0$
proves $|g_\lambda|<1$, so $\rho(G)<1$.

For every $k$, put $H_k=G^{2^k}$. Both $I-H_k$ and $I+H_k$ are
invertible, since all eigenvalues of $H_k$ have modulus less than one.
Define
\begin{equation}\label{eq:eq-sign-cayley}
 Y_k=(I-H_k)^{-1}(I+H_k)S.
\end{equation}
All factors commute. At $k=0$, the identity
$G(A+S)=A-S$ gives $(I-G)A=(I+G)S$, hence $Y_0=A$.
Each $Y_k$ is invertible, with
$Y_k^{-1}=S(I+H_k)^{-1}(I-H_k)$. Bringing these two expressions to
a common denominator gives
\[
 \tfrac12(Y_k+Y_k^{-1})
 =(I-H_k^2)^{-1}(I+H_k^2)S=Y_{k+1}.
\]
Therefore the actual Newton recurrence is defined for every step and
equals \eqref{eq:eq-sign-cayley}. The existing power-convergence theorem
gives $H_k\to0$; inversion is continuous at $I$, so $Y_k\to S$.
Finally, $Y_kS=SY_k$ and $S^2=I$ imply
\[
 Y_{k+1}-S=\tfrac12Y_k^{-1}(Y_k-S)^2,
 \qquad
 \|Y_{k+1}-S\|_2\le\tfrac12\|Y_k^{-1}\|_2\|Y_k-S\|_2^2.
\]
Since $Y_k^{-1}\to S$, these constants are eventually bounded, which
is the precise quadratic-convergence conclusion. Spectral radius below
one does not assert $\|G\|_2<1$ for a nonnormal matrix.

\paragraph{Verified reuse and source check.}
\code{NLAlib.tendsto_matrix_pow_zero_of_spectralRadius_lt_one} is
already proved for arbitrary complex square matrices. Its application
to the subsequence $2^k$ is the entire convergence input here.
The Newton recurrence is also recorded in
\href{https://nhigham.com/2020/12/15/what-is-the-matrix-sign-function/}
{Higham's matrix-sign overview}; the algebraic derivation above is
given explicitly so no matrix-function calculus is an assumed step.

\subsection{Sylvester block sign from the uniqueness characterization}
\label{sec:eq-sign-sylvester}
\textbf{Atlas id:} \code{matrix-sign-sylvester-block}.

Let $A\in\C^{m\times m}$ and $B\in\C^{n\times n}$ be positive
stable. Their spectra and the spectrum of $-B$ are disjoint, so
$AX+XB=C$ has a unique solution. Set
\[
 M=\begin{bmatrix}A&C\\0&-B\end{bmatrix},\qquad
 T=\begin{bmatrix}I&2X\\0&-I\end{bmatrix}.
\]
Block multiplication gives $T^2=I$ and
\[
 MT=\begin{bmatrix}A&2AX-C\\0&B\end{bmatrix},\qquad
 TM=\begin{bmatrix}A&C-2XB\\0&B\end{bmatrix}.
\]
They are equal precisely because $AX+XB=C$. The diagonal blocks of
$TM$ are $A$ and $B$, so $TM$ is positive stable. The spectrum of $M$
avoids the imaginary axis. The preceding uniqueness characterization
therefore proves $\operatorname{sign}(M)=T$, including the factor $2$
and its positive sign. No contour integral or block diagonalization
is needed. Empty block sizes merely delete the corresponding row and
column blocks.

\paragraph{Verified reuse.}
The required solvability theorem is
\code{NLAlib.sylvesterLinearMap_bijective_iff_disjoint_spectrum}, applied
to $(A,-B)$ over $\C$.

\subsection{Newton--Schulz inverse iteration without commutation assumptions}
\label{sec:eq-inverse-schulz}
\textbf{Atlas id:} \code{matrix-inverse-newton-schulz}.

Let $E_0=I-AX_0$ and assume $\rho(E_0)<1$. Then $1$ is not an
eigenvalue of $E_0$, so $AX_0=I-E_0$ is invertible. In finite square
dimension, invertibility of a product implies invertibility of each
factor; in particular $A$ is invertible. Define
$X_{k+1}=X_k(2I-AX_k)$ and $E_k=I-AX_k$. Associativity alone gives
\[
 E_{k+1}=I-(AX_k)(2I-AX_k)=(I-AX_k)^2=E_k^2.
\]
Induction proves $E_k=E_0^{2^k}$. Since
$X_k=A^{-1}(I-E_k)$, the power-convergence theorem yields
$X_k\to A^{-1}$. If the stronger condition $q=\|E_0\|_2<1$ holds,
submultiplicativity gives
\[
 \|X_k-A^{-1}\|_2
 \le\|A^{-1}\|_2\|E_0^{2^k}\|_2
 \le\|A^{-1}\|_2q^{2^k}.
\]
Neither $AX_0=X_0A$ nor any relation among the original factors is
needed. The finite residual identity should be formalized before the
limit theorem, since it is purely associative ring algebra.

\subsection{Pseudospectra: a rank-one correction attains the backward error}
\label{sec:eq-pseudospectrum}
\textbf{Atlas id:} \code{pseudospectrum-perturbation-equivalence}.

Let $n\ge1$, $M=zI-A$, and $s=\sigma_{\min}(M)$, with the
domain-side minimum singular value. The SVD gives a unit vector $v$
with $\|Mv\|=s$. Set
\[
 E=(Mv)v^*.
\]
Then $(A+E)v=Av+Mv=zv$, $\rank E\le1$, and
$\|E\|_2=\|Mv\|=s$: Cauchy--Schwarz gives the upper bound and
evaluation at $v$ gives equality. Conversely, if $z$ is an eigenvalue
of $A+F$, choose a unit eigenvector $w$. Then
$Mw=Fw$, and
$s\le\|Mw\|\le\|F\|_2$. Consequently
\[
 \min\{\|F\|_2:z\in\operatorname{spec}(A+F)\}=s.
\]
The same SVD minimum immediately gives
$s\le\varepsilon$ if and only if some unit $w$ has
$\|(zI-A)w\|\le\varepsilon$. This proves the residual and perturbation
descriptions of the closed $\varepsilon$-pseudospectrum, including
$\varepsilon=0$.

If $z\notin\operatorname{spec}(A)$, then $s>0$ and the inverse SVD
gives $\|(zI-A)^{-1}\|_2=1/s$. Thus, for $\varepsilon>0$, the same
condition is equivalent to
\[
 z\in\operatorname{spec}(A)
 \quad\hbox{or}\quad
 \bigl(z\notin\operatorname{spec}(A)\ \hbox{and}\
       \|(zI-A)^{-1}\|_2\ge1/\varepsilon\bigr).
\]
The explicit guard on the inverse avoids interpreting a singular
totalized inverse as an infinite resolvent norm. If $A$ is normal,
unitary diagonalization gives
$s=\min_{\lambda\in\operatorname{spec}(A)}|z-\lambda|$, proving the
union-of-closed-disks description. The perturbation class here is
unstructured complex matrices, even if $A$ happens to be real.

\paragraph{Verified reuse.}
\code{NLAlib.exists_orthonormalBasis_eigenvectors_of_normal} supplies
the final normal-matrix step. The general argument uses only the
already verified SVD and rank-one norm calculation.

\subsection{Logarithmic norm: an integrating factor, without Gronwall}
\label{sec:eq-lognorm}
\textbf{Atlas id:} \code{matrix-exponential-logarithmic-norm}.

Let $n\ge1$, $H=(A+A^*)/2$, and $\omega=\lambda_{\max}(H)$.
For a fixed vector $v$, put $y(t)=e^{tA}v$. The exponential derivative
gives $y'(t)=Ay(t)$, and differentiation of the real scalar squared
norm yields
\[
 \frac{d}{dt}\|y(t)\|^2
 =2\operatorname{Re}(y(t)^*Ay(t))
 =2y(t)^*Hy(t)\le2\omega\|y(t)\|^2.
\]
Consequently the derivative of
$e^{-2\omega t}\|y(t)\|^2$ is nonpositive. The ordinary scalar
mean-value theorem makes this function nonincreasing, so
$\|e^{tA}v\|\le e^{\omega t}\|v\|$ for all $t\ge0$.
Taking the operator norm proves $\|e^{tA}\|_2\le e^{\omega t}$.

To prove optimality of the exponent, choose a unit eigenvector $v$ of
$H$ for $\omega$. Suppose $\|e^{tA}\|_2\le e^{ct}$ for every
$t\ge0$. Then $f(t)=\|e^{tA}v\|^2\le e^{2ct}$, with equality at
$t=0$. Subtract that common value, divide by $t>0$, and let
$t\downarrow0$. Since $f'(0)=2\omega$, this gives $\omega\le c$.
Thus $\omega$ is exactly the minimum admissible exponent with leading
constant one. Applying the same result with $c=0$ proves
\[
 \bigl(\forall t\ge0,\ \|e^{tA}\|_2\le1\bigr)
 \quad\Longleftrightarrow\quad H\preceq0.
\]
The derivative at zero supplies necessity; a spectral-abscissa estimate
would not suffice for a nonnormal matrix.

\paragraph{Verified reuse and source check.}
Mathlib's
\code{hasDerivAt_exp_smul_const} and
\code{hasDerivAt_exp_smul_const'} are present in
\code{Mathlib.Analysis.SpecialFunctions.Exponential}; their hypotheses
allow noncommutative Banach algebras. They specialize directly to real
time and complex matrices. The sharp growth statement and the Hermitian
formula agree with Theorems 3 and 6 of
\href{https://nhigham.com/2022/01/18/what-is-the-logarithmic-norm/}
{Higham's logarithmic-norm overview}. Only one scalar derivative
inequality is needed in the derivation above.

\section{Pencils and constrained systems by finite block algebra}
\label{sec:eq-blocks}

\subsection{Complex QZ: one dependent image pair at each induction step}
\label{sec:eq-qz}
\textbf{Atlas id:} \code{generalized-schur-qz}.

We first prove a small lemma for linear maps $A,B:E\to F$ between
nonzero complex spaces of the same finite dimension: there are unit
vectors $v\in E$, $u\in F$ such that both $Av$ and $Bv$ lie in the
line spanned by $u$.
If $B$ is singular, take a unit $v\in\ker B$. When $Av\ne0$, take
$u=Av/\|Av\|$; when $Av=0$, take any unit $u$. If $B$ is invertible,
the endomorphism $B^{-1}A$ has a nonzero eigenvector $v$ over $\C$.
Normalize it, so $Av=\lambda Bv$, and normalize the nonzero vector
$Bv$ to obtain $u$. This proves the lemma with no regular-pencil
assumption.

Now use induction on $n$. For $n=0$ use the empty bases. For $n>0$,
choose $u,v$ as above and orthogonally decompose
$E=\mathbb Cv\oplus v^\perp$, $F=\mathbb Cu\oplus u^\perp$.
The two matrices in these decompositions have the shapes
\[
 \begin{bmatrix}a&*\\0&A_1\end{bmatrix},\qquad
 \begin{bmatrix}b&*\\0&B_1\end{bmatrix},
\]
where $A_1=P_{u^\perp}A|_{v^\perp}$ and
$B_1=P_{u^\perp}B|_{v^\perp}$. The compressed maps again have equal
domain and codomain dimension. Apply the induction hypothesis to
$(A_1,B_1)$ and extend the resulting orthonormal bases by $v,u$.
The resulting unitary matrices $Q,Z$ satisfy
$S=Q^*AZ$, $T=Q^*BZ$, both upper triangular.

For any scalars $\alpha,\beta$, determinant multiplicativity gives
\begin{equation}\label{eq:eq-qz-det}
 \det(\beta A-\alpha B)
 =\det(Q)\det(Z^*)\prod_{i=1}^n(\beta s_{ii}-\alpha t_{ii}).
\end{equation}
If $\det(A-\lambda B)$ is not the zero polynomial, none of the pairs
$(s_{ii},t_{ii})$ is $(0,0)$, because such a pair would make the product
identically zero. Conversely, if no pair is zero, each factor
$s_{ii}-\lambda t_{ii}$ is a nonzero polynomial and their product is
nonzero. Thus regularity is exactly the condition excluding a zero
diagonal pair. The homogeneous factorization
\eqref{eq:eq-qz-det} gives the projective eigenvalues
$[s_{ii}:t_{ii}]$, including infinity when $t_{ii}=0$, with
multiplicities. The dimension-zero determinant is $1$ and the list of
projective eigenvalues is empty. Singular pencils still admit the
triangularization, but the projective-root interpretation is then not
asserted.

\paragraph{Verified reuse: the induction already exists.}
The live declaration index contains, all with \code{sorry=false},
\code{NLAlib.exists_nonzero_pencil_image_line},
\code{NLAlib.exists_unit_pencil_image_line}, and
\code{NLAlib.exists_orthonormalBases_pencil_upper_triangular}.
The last theorem is already stated for arbitrary pairs of equal
finite-dimensional complex spaces. Thus this open card primarily needs
the matrix wrapper, determinant factorization, and explicit regularity
interpretation; its central induction should be reused.

\subsection{KKT systems: identify the kernel using only the null space}
\label{sec:eq-kkt}
\textbf{Atlas id:} \code{kkt-reduced-hessian}.

Let $H=H^*$ and let $B:\mathbb K^n\to\mathbb K^m$ have full row
rank. Let $Z$ have orthonormal columns spanning $\ker B$, and put
$G=Z^*HZ$. The map $B^*$ is injective and
\[
 (\ker B)^\perp=\operatorname{range}B^*,\qquad BB^*\text{ is invertible}.
\]
For $K=\begin{bmatrix}H&B^*\\B&0\end{bmatrix}$, a vector in its
kernel obeys $Bx=0$, hence $x=Zy$, and
$Hx+B^*\lambda=0$. Multiplying by $Z^*$ gives $Gy=0$.
If $G$ is invertible then $y=x=0$, and injectivity of $B^*$ gives
$\lambda=0$, so $K$ is invertible.
Conversely, if $Gy=0$ for nonzero $y$, then $x=Zy\ne0$ and
$Z^*Hx=0$, so $Hx\in\operatorname{range}B^*$. Choosing
$\lambda=-(BB^*)^{-1}BHx$ gives $Hx+B^*\lambda=0$.
Thus $K$ is singular. This proves the exact invertibility equivalence.

For a right-hand side $(f,g)$, full row rank supplies a feasible
$x_p$, for example $x_p=B^*(BB^*)^{-1}g$. When $G$ is invertible set
\[
 y=G^{-1}Z^*(f-Hx_p),\qquad
 x=x_p+Zy,\qquad
 \lambda=(BB^*)^{-1}B(f-Hx).
\]
Then $Bx=g$ and $Z^*(f-Hx)=0$, so $f-Hx\in\operatorname{range}B^*$.
The multiplier formula is therefore exact and yields
$Hx+B^*\lambda=f$.

For the real-valued objective
$J(x)=\tfrac12x^*Hx-\operatorname{Re}(f^*x)$, restricted to $Bx=g$,
expansion at the above solution gives
\[
 J(x+Zz)-J(x)=\tfrac12 z^*Gz.
\]
Hence $G\succ0$ implies that $x$ is the unique constrained minimum.
Conversely, if the problem has a unique minimum for every $f,g$, take
$f=g=0$. A negative eigenvalue of $G$ makes the objective unbounded
below along its eigenvector. If $G\succeq0$ is singular, a nonzero
kernel vector gives multiple minimizers. Thus $G\succ0$ is necessary
as well. When $n=m$, $Z$ has zero columns: $G$ is the unique operator
on the zero space, both invertible and vacuously positive definite,
and the feasible set is a singleton. When $m=0$, the result reduces to
the unconstrained Hermitian problem.

\subsection{Sylvester separation: the least gain of one finite operator}
\label{sec:eq-separation}
\textbf{Atlas id:} \code{sylvester-separation-residual-conditioning}.

Let $m,n\ge1$ and define the linear map
$\mathcal L(Z)=AZ-ZB$ on the $mn$-dimensional space of $m\times n$
matrices with Frobenius norm. Define
\[
 \delta=\operatorname{sep}_F(A,B)
 =\min_{\|Z\|_F=1}\|\mathcal L(Z)\|_F.
\]
The minimum exists by compactness of a finite-dimensional unit sphere.
Equivalently it is the smallest singular value of $\mathcal L$ in
the orthonormal matrix-unit basis. Column-by-column vectorization is
an isometry, and entrywise multiplication verifies
\[
 \operatorname{vec}(AZ-ZB)
 =(I_n\otimes A-B^T\otimes I_m)\operatorname{vec}(Z).
\]
This proves the Kronecker formula, with a transpose, not an adjoint.
Homogeneity gives $\|\mathcal L(Z)\|_F\ge\delta\|Z\|_F$ for all $Z$.
Also $\delta>0$ if and only if $\mathcal L$ is injective: if the
minimum is zero, its attaining unit vector is in the kernel; the
converse is immediate. Finite equal dimensions make this equivalent
to invertibility.

Over $\C$, the existing Sylvester theorem identifies invertibility
with disjoint spectra. One elementary verification is to triangularize
$A$ and $B$ separately; vectorization produces a triangular matrix,
after reversing one of the lexicographic orders if needed, with
diagonal entries $\lambda_i(A)-\lambda_j(B)$. Its determinant is their
product. If spectra meet, one can also directly use a right eigenvector
$u$ of $A$ and a left eigenvector $w^*$ of $B$ for the same eigenvalue:
$uw^*$ is a nonzero kernel matrix.

\paragraph{Necessary real-spectrum clarification.}
For real $A,B$, the spectral assertion must use their spectra after
complexification. Indeed, complex and real injectivity are equivalent:
a complex kernel matrix has real and imaginary parts in the real
kernel, and at least one part is nonzero. The least gain is unchanged
as well: for $Z=X+iY$, the real coefficients give
\[
 \|\mathcal L(Z)\|_F^2
 =\|\mathcal L(X)\|_F^2+\|\mathcal L(Y)\|_F^2
 \ge\delta_{\R}^2(\|X\|_F^2+\|Y\|_F^2),
\]
while real vectors give the reverse inequality of the minima.
Using only real eigenvalues would be false: take
$A=B=\begin{bmatrix}0&-1\\1&0\end{bmatrix}$. Both real spectra are
empty, but $\mathcal L(I)=0$ and $\delta=0$.

Now assume $\delta>0$ and $\mathcal L(X)=C$. If
$R=C-AY+YB$, then $R=\mathcal L(X-Y)$ and
\begin{equation}\label{eq:eq-sep-residual}
 \|X-Y\|_F\le\frac{\|R\|_F}{\delta}.
\end{equation}
The constant is sharp: choose a unit matrix $Z_0$ attaining the minimum
and take $C=\mathcal L(Z_0)$, $X=Z_0$, $Y=0$. Equivalently
$\|\mathcal L^{-1}\|_{F\to F}=1/\delta$.

For perturbations define
$\Delta\mathcal L(Z)=\Delta A\,Z-Z\Delta B$ and
$d=\|\Delta A\|_2+\|\Delta B\|_2$. The mixed norm inequalities give
$\|\Delta\mathcal L(Z)\|_F\le d\|Z\|_F$, and hence
\[
 \|(\mathcal L+\Delta\mathcal L)(Z)\|_F
 \ge(\delta-d)\|Z\|_F.
\]
If $d<\delta$, the perturbed operator is injective and therefore
invertible; no infinite Neumann series is required. If
$(\mathcal L+\Delta\mathcal L)(\widehat X)=C+\Delta C$, subtraction
gives
\[
 (\mathcal L+\Delta\mathcal L)(\widehat X-X)
 =\Delta C-\Delta\mathcal L(X).
\]
Applying its lower-gain estimate proves
\[
 \|\widehat X-X\|_F
 \le\frac{\|\Delta C\|_F+d\|X\|_F}{\delta-d}.
\]
When $A,B$ are normal, their separate unitary eigenbases turn
$\mathcal L$ into entrywise multiplication by $\lambda_i-\mu_j$;
therefore $\delta=\min_{i,j}|\lambda_i-\mu_j|$. This equality is not
asserted for nonnormal matrices. The residual norm here is Frobenius;
the same constant has not been asserted for the spectral norm of a
matrix residual.

\paragraph{Verified reuse.}
Use \code{NLAlib.sylvesterLinearMap},
\code{NLAlib.sylvesterLinearMap_bijective_iff_disjoint_spectrum},
\code{NLAlib.frobNorm_mul_le_specNorm_mul_frobNorm}, and
\code{NLAlib.frobNorm_mul_le_frobNorm_mul_specNorm}.
The bijectivity theorem assumes an algebraically closed field, so the
real-complex bridge just proved is a genuine required interface, not
an implicit application over $\R$.



\section{Nonnegative matrices: one resolvent construction and finite certificates}
\label{sec:ns-perron}

This section covers \code{nonnegative-perron-eigenpair},
\code{irreducible-perron-frobenius},
\code{nonnegative-spectral-radius-comparison},
\code{collatz-wielandt-certificates}, and
\code{primitive-perron-power-limit}. All inequalities between real matrices
or vectors in this section are entrywise. The dimension is positive, and
$\rho(A)$ means the spectral radius of the complexification. A directed
edge $j\to i$ means $A_{ij}>0$; reversing all edges makes no difference to
strong connectivity.

\subsection{A normalized resolvent proves Perron existence, including radius zero}

\begin{theorem}[Nonnegative Perron pair]
If $A\in\R^{n\times n}$, $n\ge1$, and $A\ge0$, there is $p\ge0$ such that
$\sum_i p_i=1$ and $Ap=rp$, where $r=\rho(A)$. The same assertion for
$A^\top$ supplies a nonnegative left Perron vector.
\end{theorem}
\begin{proof}
For $t>r$ the existing spectral-radius Neumann theorem gives
\[
 R_t=(tI-A)^{-1}=\sum_{k\ge0}t^{-k-1}A^k\ge0,
 \qquad v_t=R_t\boldsymbol1\ge t^{-1}\boldsymbol1>0.
\]
The missing point is to prove that $v_t$ becomes unbounded without already
knowing that $r$ is an eigenvalue. Choose a complex eigenpair
$Az=\lambda z$, $z\ne0$, with $|\lambda|=r$, and put
$M=\max_i|z_i|>0$. Positivity and the triangle inequality give, for every
$k\ge0$,
\[
 A^k\boldsymbol1\ge M^{-1}A^k|z|
       \ge M^{-1}|A^kz|=M^{-1}r^k|z|.
\]
At a coordinate with $|z_i|=M$, summing yields
$(v_t)_i\ge1/(t-r)$. This includes $r=0$; the $k=0$ term alone then gives
the required bound. Consequently $s_t=\sum_i(v_t)_i\ge1/(t-r)$.

Take $t_j=r+1/(j+1)$ and $p_j=v_{t_j}/s_{t_j}$. The finite simplex is
compact, so a subsequence converges to a vector $p$ in that simplex. The
identity
\[
 A p_j=t_jp_j-s_{t_j}^{-1}\boldsymbol1
\]
passes to the limit and gives $Ap=rp$. No fixed-point theorem, perturbation
of roots, or irreducibility assumption is used.
\end{proof}

\paragraph{Verified reuse.}
The snapshot contains the proved declarations
\code{NLAlib.hasSum_real_matrix_pow_inv_of_complex_spectralRadius_lt_one},
\code{NLAlib.summable_real_matrix_pow_of_complex_spectralRadius_lt_one},
and \code{NLAlib.tendsto_real_matrix_pow_zero_iff_complex_spectralRadius_lt_one}.
The new content is the displayed entrywise bound and compactness of a finite
simplex. Positivity of a convergent series follows by taking limits of its
nonnegative partial sums.

\subsection{Irreducibility, algebraic simplicity, and the projector}

Suppose now that $A$ is irreducible and nonzero. If $v\ge0$, $v\ne0$,
and $Av=\lambda v$, a zero coordinate forces every predecessor along a
positive entry to have zero coordinate. Strong connectivity would force
$v=0$. Thus every such eigenvector is strictly positive. Applying this to
the right and left Perron vectors gives $p,q>0$. Since $A\ne0$, $Ap$ has
a positive coordinate, and therefore $r>0$.

For a real eigenvector $u$ at $r$, let $a=\max_i u_i/p_i$. Then
$ap-u\ge0$, it has a zero coordinate, and it is an eigenvector at $r$.
The previous paragraph forces $ap-u=0$. Real and imaginary parts extend
this statement to complex eigenvectors:
\[
 \ker(A-rI)=\C p.
\]
Furthermore, $q^\top(A-rI)=0$ and $q^\top p>0$, so $p$ does not belong
to $\operatorname{range}(A-rI)$. Rank--nullity now gives the invariant
direct sum
\[
 \C^n=\C p\oplus\operatorname{range}(A-rI).
\]
The restriction of $A-rI$ to its range has zero kernel, hence is invertible.
In a basis adapted to this direct sum, the characteristic polynomial has
exactly one factor $X-r$. This proves \emph{algebraic}, not just geometric,
simplicity without Jordan normal form.

The invariant complementary space is also $\ker q^\top$, since both have
dimension $n-1$. Thus the spectral projector at $r$ is
\[
 P=\frac{pq^\top}{q^\top p},\qquad P^2=P,\qquad AP=PA=rP,
 \qquad P_{ij}>0.
\]
This projector need not be orthogonal. If $Aw=\lambda w$ with
$0\ne w\ge0$, multiplication by $q^\top$ gives $\lambda=r$; the
maximum-ratio argument then gives $w=ap$ with $a>0$. The left-vector
statement follows by transposition. The explicit nonzero hypothesis
excludes the one-dimensional zero matrix; irreducibility alone does not
exclude it under the usual graph convention.

\subsection{Comparison and Collatz--Wielandt without limits of spectra}

For $t>\rho(A)$ use the strictly positive vector
$v=(tI-A)^{-1}\boldsymbol1$, so $Av=tv-\boldsymbol1<tv$.
If $|B_{ij}|\le A_{ij}$ and $Bz=\lambda z$, take an index attaining
$c=\max_i|z_i|/v_i>0$. At that index,
\[
 |\lambda||z_i|\le(A|z|)_i\le c(Av)_i<t c v_i=t|z_i|.
\]
Hence $|\lambda|<t$ for every $t>\rho(A)$ and
$\rho(B)\le\rho(A)$. This proves entrywise monotonicity as a special case.
For strict monotonicity, let $A$ be irreducible and nonzero, $E\ge0$,
$E\ne0$, and set $C=A+E$. Then $C$ is irreducible. Use a positive right
Perron vector $p$ for $A$ and a positive left Perron vector $q_C$ for $C$:
\[
 (\rho(C)-\rho(A))q_C^\top p=q_C^\top E p>0.
\]

For $x>0$ set
$\alpha(x)=\min_i(Ax)_i/x_i$ and
$\beta(x)=\max_i(Ax)_i/x_i$. A nonnegative nonzero left Perron vector
$q$ satisfies $q^\top x>0$; multiplying
$\alpha(x)x\le Ax\le\beta(x)x$ by $q^\top$ proves
\[
 \alpha(x)\le\rho(A)\le\beta(x).
\]
The resolvent vector above has $\beta(v)<t$. Letting $t\downarrow\rho(A)$
in this scalar inequality proves
\[
 \rho(A)=\inf_{x>0}\beta(x).
\]
If $A$ is irreducible and nonzero, $x=p>0$ attains both bounds, so
$\min_{x>0}\beta(x)=\max_{x>0}\alpha(x)=\rho(A)$.

For the general support-aware formula, take $0\ne x\ge0$ and put
$a=\min_{i:x_i>0}(Ax)_i/x_i$. Then $Ax\ge ax$ at every coordinate,
including those outside the support. For $v=(tI-A)^{-1}\boldsymbol1$,
choose an index attaining $c=\max_i x_i/v_i>0$. At this index,
\[
 a x_i\le(Ax)_i\le c(Av)_i<t x_i.
\]
Thus $a\le\rho(A)$. Taking the possibly non-strictly-positive Perron
vector $x=p$ attains equality, proving
\[
 \rho(A)=\max_{0\ne x\ge0}\min_{i:x_i>0}\frac{(Ax)_i}{x_i}.
\]
The support qualification is essential: for $A=\diag(0,1)$,
$\alpha(x)=0$ for every $x>0$, whereas $\rho(A)=1$.

\subsection{Primitivity and the power limit through a stochastic similarity}

Continue with irreducible nonzero $A\ge0$, and normalize
$\sum_i p_i=1$. Put
\[
 T=A/r,\qquad D=\diag(p),\qquad K=D^{-1}TD.
\]
The matrix $K$ is nonnegative and its row sums are one. If some $A^m>0$,
then $K^m>0$. Suppose $Kz=\lambda z$ with $|\lambda|=1$. Normalize
$\max_i|z_i|=1$ and choose a maximizing coordinate $i$.
Writing $a_j=(K^m)_{ij}>0$ gives
$|\sum_j a_jz_j|=1$, $\sum_j a_j=1$. The finite identity
\[
 \sum_j a_j|z_j-c|^2=\sum_j a_j|z_j|^2-|c|^2,
 \qquad c=\sum_j a_jz_j,
\]
forces $z_j=c$ for every $j$. Hence $z$ is constant and $\lambda=1$.
Primitivity therefore excludes all other peripheral eigenvalues.

Conversely, suppose $1$ is the only eigenvalue of $T$ of modulus one.
The direct sum $\C p\oplus\ker q^\top$ above shows that, with
$P=pq^\top/(q^\top p)$ and $Q=T-P$, all eigenvalues of $Q$ have modulus
strictly less than one. The existing power-convergence theorem yields
$Q^k\to0$. Since $TP=PT=P$,
\[
 T^k=P+Q^k\quad(k\ge1),\qquad T^k\longrightarrow P>0.
\]
Entrywise convergence in a finite matrix implies that some $T^k$ is
entrywise positive. Thus the spectral-circle condition implies primitivity.
Finally, if $T^k$ merely converges, then for every peripheral eigenpair
$Tz=\lambda z$, the sequence $\lambda^kz$ converges; its nonzero norm
forces $\lambda=1$ by taking the limit in $\lambda^{k+1}=\lambda\lambda^k$.
This proves all three equivalences in the card.

For $0\ne x\ge0$, no power $A^kx$ vanishes: $q^\top A^kx=r^kq^\top x>0$.
Moreover
$T^kx\to p(q^\top x)/(q^\top p)$, whose coordinate sum is positive.
Continuity of scalar division gives
\[
 \frac{A^kx}{\sum_i(A^kx)_i}\longrightarrow p.
\]
For an explicit admissible rate, fix $\rho(Q)<\theta<1$ and choose $N$
such that $\norm{Q^k}_2\le\theta^k$ for all $k\ge N$; the Gelfand
formula provides this $N$. Define the finite constant
\[
 C_\theta=\max\left(1,\norm{I-P}_2,
             \max_{0\le k<N}\theta^{-k}\norm{Q^k}_2\right),
\]
where the inner maximum is omitted if $N=0$. Then
$\norm{T^k-P}_2\le C_\theta\theta^k$ for every $k\ge0$.
The separate $k=0$ term avoids the false identity $T^0-P=Q^0$.
Verified reuse includes
\code{NLAlib.tendsto_norm_matrix_pow_rpow_spectralRadius},
\code{NLAlib.exists_prefix_norm_matrix_pow_le_spectralRadius_add}, and
\code{NLAlib.norm_matrix_pow_le_matrixPowerBoundConstant_mul_pow}.
No assertion of convergence is made from irreducibility alone.

\section{Nonsingular M-matrices from positive scaling and scalar elimination}
\label{sec:ns-mmatrix}

This covers \code{nonsingular-m-matrix-characterizations} and
\code{m-matrix-schur-complement}. Let $A\in\R^{n\times n}$, $n\ge1$,
have nonpositive off-diagonal entries. Choose, once and for all,
$s=1+\max(0,\max_i A_{ii})$ and $B=sI-A\ge0$.

\begin{theorem}[Five equivalent certificates]
The following conditions are equivalent:
\begin{enumerate}
\item $A$ is invertible and $A^{-1}\ge0$;
\item $A=s'I-B'$ for some $B'\ge0$ and $s'>\rho(B')$;
\item every complex eigenvalue of $A$ has positive real part;
\item every nonempty principal minor of $A$ is positive;
\item some $x>0$ satisfies $Ax>0$.
\end{enumerate}
They are also equivalent to strict row diagonal dominance with positive
diagonal of $A\diag(x)$ for some $x>0$.
\end{theorem}
\begin{proof}
If $A^{-1}\ge0$, each row of the inverse is nonzero, so
$x=A^{-1}\boldsymbol1>0$ and $Ax=\boldsymbol1$.
Conversely, $Ax>0$ means $Bx<sx$; the positive-vector upper certificate
gives $\rho(B)<s$. The Neumann series then proves invertibility and
\[
 A^{-1}=s^{-1}\sum_{k\ge0}(B/s)^k\ge0.
\]
This establishes the equivalence of conditions 1, 2 and 5.
For any such representation, $s'-\lambda$ is an eigenvalue of $A$
whenever $\lambda$ is an eigenvalue of $B'$, so its real part is at least
$s'-\rho(B')>0$. In the other direction use the fixed representation
$A=sI-B$: the Perron root $\rho(B)$ is an eigenvalue of $B$, hence
$s-\rho(B)>0$ by condition 3.

To obtain principal minors from condition 2, every principal block
$B'[J,J]$ has spectral radius at most $\rho(B')$: embed it with zeros in
the full matrix and apply entrywise domination. For $0\le t\le1$,
$s'I-tB'[J,J]$ is invertible. Its real determinant is continuous, never
zero, and positive at $t=0$. Therefore
$\det A[J,J]>0$.

For the reverse implication from condition 4, induct on $n$. The
one-dimensional case is immediate. Write
\[
 A=\begin{pmatrix}a&f\\g&H\end{pmatrix},\qquad a>0,
 \qquad S=H-ga^{-1}f.
\]
The off-diagonal entries of $S$ are nonpositive, since $f,g\le0$.
For each nonempty principal index set $J$ in the trailing block,
\[
 \det S[J,J]=\frac{\det A[\{1\}\cup J,\{1\}\cup J]}{a}>0.
\]
Induction gives $S^{-1}\ge0$. The elementary block inverse is
\[
 A^{-1}=\begin{pmatrix}
 a^{-1}+a^{-1}fS^{-1}ga^{-1}&-a^{-1}fS^{-1}\\
 -S^{-1}ga^{-1}&S^{-1}
 \end{pmatrix}\ge0.
\]
This induction uses only scalar pivots and determinant row elimination.
Finally, because $A_{ij}\le0$ for $i\ne j$,
\[
 (Ax)_i>0
 \quad\Longleftrightarrow\quad
 A_{ii}x_i>\sum_{j\ne i}|A_{ij}|x_j.
\]
The right side implies $A_{ii}>0$ and is precisely the claimed scaling
condition.
\end{proof}

Multiplying by the nonnegative inverse proves the monotonicity statement
$Au\ge0\Rightarrow u\ge0$. No symmetry or Loewner-order argument is
needed.

For the block result write $A=\left(\begin{smallmatrix}E&F\\G&H\end{smallmatrix}\right)$
with nonempty principal blocks and $S=H-GE^{-1}F$. If $A$ is a nonsingular
M-matrix, every principal block is one by the principal-minor criterion.
Thus $E^{-1}\ge0$, and $S$ is a Z-matrix. The block inverse identifies
$S^{-1}$ with the lower-right block of $A^{-1}$, so $S^{-1}\ge0$ and $S$
is a nonsingular M-matrix. Conversely, if $E$ and $S$ are nonsingular
M-matrices, the same block inverse formula gives
\[
 A^{-1}=\begin{pmatrix}
 E^{-1}+E^{-1}FS^{-1}GE^{-1}&-E^{-1}FS^{-1}\\
 -S^{-1}GE^{-1}&S^{-1}
 \end{pmatrix}\ge0.
\]
This proves the equivalence and the inverse-block identity. It also
clarifies the hypothesis order: define $S$ only after assuming $E$ is
invertible. Empty eliminated blocks, if added as variants, simply return
the original matrix and require no new argument.

\section{Fourier structure: one character sum, scalar division, and index splits}
\label{sec:ns-fourier}

This section covers \code{circulant-fourier-diagonalization},
\code{circulant-fourier-pseudoinverse},
\code{radix-two-fft-realization}, and
\code{toeplitz-circulant-fft-product}.

\subsection{Diagonalization and all spectral consequences}

Let $n\ge1$, $\omega=e^{-2\pi i/n}$,
$F_{jk}=\omega^{jk}$, $U=F/\sqrt n$, and
$C_{ij}=c_{i-j}$ on $\mathbb Z/n\mathbb Z$. The single finite identity
\[
 \sum_{j=0}^{n-1}\omega^{aj}=
 \begin{cases}n&a=0\pmod n,\\0&a\ne0\pmod n\end{cases}
\]
follows from a geometric sum and proves $FF^*=nI$, hence $U$ is unitary.
Changing the variable $h=i-j$ gives directly
\[
 (FC)_{kj}=\sum_i\omega^{ki}c_{i-j}
       =\omega^{kj}\sum_h\omega^{kh}c_h
       =(Fc)_kF_{kj}.
\]
Therefore, with $\lambda=Fc$,
\[
 UCU^*=\diag(\lambda),\qquad
 UC^*CU^*=\diag(|\lambda_j|^2).
\]
The singular-value multiset is $\{ |\lambda_j|\}_j$, with repetitions
retained. The Rayleigh quotient of the last diagonal matrix gives
$\norm C_2=\max_j|\lambda_j|$. Invertibility is equivalent to all
$\lambda_j\ne0$; then
$\norm{C^{-1}}_2=(\min_j|\lambda_j|)^{-1}$ and
$\kappa_2(C)=\max_j|\lambda_j|/\min_j|\lambda_j|$.
These statements include $n=1$ and repeated modes. The zero-dimensional
case is excluded by the Fourier convention and by the maxima/minima.

\paragraph{Verified reuse and conventions.}
Mathlib's \code{ZMod.dft}, \code{ZMod.dft_apply},
\code{ZMod.invDFT_apply}, \code{ZMod.dft_dft}, and
\code{Matrix.circulant} were checked in the local source. The DFT uses
the negative exponent and no normalization; the circulant entry is
$c_{i-j}$. Thus the diagonal is $Fc$, not $Uc$. The matrix identity is a
bridge to those existing definitions, not a reason to define another DFT.

\subsection{Pseudoinverse and least squares are independent scalar problems}

Define $d_j=1/\lambda_j$ when $\lambda_j\ne0$ and $d_j=0$ otherwise.
The scalar diagonal matrices $\diag(\lambda)$ and $\diag(d)$ satisfy
the four Penrose equations coordinatewise. Conjugating by $U$ and using
uniqueness therefore gives
\[
 C^+=U^*\diag(d)U=\operatorname{circulant}(F^{-1}d).
\]
The second equality follows by the just-proved diagonalization applied
to the first column $F^{-1}d$. For $\eta=Ub$ and $\xi=Ux$, unitary
invariance gives
\[
 \norm{Cx-b}_2^2=\sum_j|\lambda_j\xi_j-\eta_j|^2.
\]
At nonzero modes the minimizing coordinate is uniquely
$\xi_j=d_j\eta_j$; at zero modes the residual is fixed and minimum
solution norm requires $\xi_j=0$. Consequently
\[
 x_*=U^*\diag(d)Ub=F^{-1}(d\odot Fb),\qquad
 \norm{Cx_*-b}_2^2=\sum_{\lambda_j=0}|(Ub)_j|^2.
\]
Solvability is equivalent to $(Fb)_j=0$ on all zero modes. If there are
no zero modes, the displayed solution is the unique exact solution.
The zero circulant is included: its minimum-norm least-squares solution
is zero and its residual norm is $\norm b_2$.
Verified reusable declarations are
\code{NLAlib.IsMoorePenroseInverse.unitary_transform} and
\code{NLAlib.IsMoorePenroseInverse.unique}, both over an arbitrary
\code{RCLike} field. The direct scalar least-squares proof avoids relying
on a real-only least-squares interface.

\subsection{Radix two: specification first, cost by a scalar recurrence}

For $n=2m$, split the defining DFT sum into even and odd indices. Since
$\omega_n^2=\omega_m$ and $\omega_n^m=-1$,
\[
 (F_nx)_k=(F_mx_{\mathrm{even}})_k
            +\omega_n^k(F_mx_{\mathrm{odd}})_k,
\qquad
 (F_nx)_{k+m}=(F_mx_{\mathrm{even}})_k
            -\omega_n^k(F_mx_{\mathrm{odd}})_k
\]
for $0\le k<m$. Induction on $s$ proves that the stated recurrence at
length $n=2^s$ computes exactly $F_nx$; the $s=0$ case is the identity.
This proof needs only a finite even--odd indexing equivalence and the two
root identities. It need not reason about an in-place array program.

With roots supplied as constants, each butterfly uses one multiplication
and two additions/subtractions. Thus
\[
 M_0=A_0=0,\qquad M_s=2M_{s-1}+2^{s-1},
 \qquad A_s=2A_{s-1}+2^s,
\]
and induction gives $M_s=s2^{s-1}$ for $s\ge1$ and $A_s=s2^s$.
These are exact counts for this recurrence and upper bounds when trivial
root multiplications are omitted. Root generation, data permutation,
conjugation, and bit complexity are not included in this cost model.

\subsection{Rectangular Toeplitz products: no wraparound by an interval bound}

Let $m,n\ge1$, $T_{ij}=t_{j-i}$, and $L\ge m+n-1$. Define the
first column $c$ by
\[
 c_0=t_0,\quad c_k=t_{-k}\ (1\le k<m),\quad
 c_{L-j}=t_j\ (1\le j<n),
\]
and set all other coordinates to zero. The two nonzero index ranges are
disjoint because $L-n+1\ge m$. For $0\le i<m$, $0\le j<n$,
the residue of $i-j$ is $i-j$ when $i\ge j$ and $L-(j-i)$ otherwise.
Both cases give $c_{i-j\bmod L}=t_{j-i}$. If $E_n$ pads with zeros and
$P_m$ crops, this proves entrywise
\[
 T=P_m\operatorname{circulant}(c)E_n,\qquad
 Tx=P_mF_L^{-1}\big((F_Lc)\odot(F_LE_nx)\big).
\]
For $L=2^s$, each transform has the counts above. An inverse transform
is a conjugated forward transform followed by $L$ scalar rescalings;
with $F_Lc$ precomputed, the product uses two transforms, $L$ pointwise
multiplications, and at most $L$ additional rescalings. Thus its scalar
arithmetic count is at most $3Ls+2L$ for $s\ge1$ in this straightforward
implementation. This supplies the advertised $O(L\log L)$ consequence
with its counting conventions explicit. For $L=1$, use the scalar product
directly. Cropping a circulant inverse does not follow from this identity.

\section{Hierarchical matrices: finite partitions, projections, and kernel bijections}
\label{sec:ns-hierarchical}

This section covers \code{hodlr-storage-matvec},
\code{hodlr-levelwise-compression-error},
\code{hss-nested-bases-matvec},
\code{block-separable-inverse-reduction}, and
\code{hss-inverse-rank-preservation}. The field is $\R$ or $\C$ whenever
adjoints or norms are used; the purely algebraic inverse-rank arguments
work over any field. A cluster tree is a finite rooted full binary tree.
Its root indexes all $N$ coordinates, each internal cluster is the
disjoint union of its children, and its leaves partition the root.

\subsection{HODLR: the least-common-ancestor partition is the whole proof}

For each leaf $\tau$ store the dense diagonal block
$D_\tau=A[I_\tau,I_\tau]$. For siblings $\alpha,\beta$ store exact
factorizations
\[
 A[I_\alpha,I_\beta]=U_{\alpha\beta}V_{\alpha\beta}^*,\qquad
 A[I_\beta,I_\alpha]=U_{\beta\alpha}V_{\beta\alpha}^*,
\]
with at most $k$ columns in each pair. Every ordered pair of coordinates
belongs to exactly one stored block: either both are in the same leaf,
or their least common ancestor places them in its two different children.
Consequently the sum of the embedded leaf blocks and the embedded
directed sibling products is exactly $A$. This also proves uniqueness
of the reconstructed matrix, though its factors themselves need not be
unique.

Applying each leaf block and each product $U(V^*x)$ and accumulating on
the corresponding rows therefore returns $Ax$. The stored-entry count
is at most
\[
 S=\sum_{\tau\ {\rm leaf}}|I_\tau|^2
       +\sum_{\tau\ {\rm internal}}2k|I_\tau|.
\]
Indeed one directed block stores $k(|I_\alpha|+|I_\beta|)$ scalars,
and there are two directions. A dense $a\times b$ matrix-vector action
accumulated into a vector uses at most $2ab$ scalar multiplications and
additions; zero-column actions are omitted. Applying this rule to every
stored factor gives at most $2S$ operations. If leaves have size at most
$b$ and internal depths lie in $\{0,\ldots,L-1\}$, the clusters at each
fixed depth are disjoint, so
\[
 \sum_{\rm leaf}|I_\tau|^2\le bN,\qquad
 \sum_{\rm internal}|I_\tau|\le NL,\qquad
 S\le Nb+2NkL.
\]
This proof uses the actual depth $L$. It assigns no logarithmic depth to
an unbalanced tree. It includes zero ranks and a tree whose root is
already a leaf.

\subsection{Compression errors: direct sums at a level}

Let $E=A-\widehat A$ use the same block partition. Write
$E=E_D+\sum_{\ell=0}^{L-1}E_\ell$, where $E_D$ contains the leaf
diagonal errors and $E_\ell$ the sibling errors at depth $\ell$.
The elementary two-block identity
\[
 \left\|\begin{pmatrix}0&X\\Y&0\end{pmatrix}\right\|_2
       =\max(\norm X_2,\norm Y_2)
\]
follows by applying the matrix to $(u,v)$ and using
$\norm{Xv}_2^2+\norm{Yu}_2^2$; equality follows by restricting to one
coordinate block and taking its operator-norm supremum. This includes
zero blocks. A block-diagonal direct sum likewise has norm the maximum
of its block norms. Thus, if every directed error at depth $\ell$ has
norm at most $\varepsilon_\ell$ and every leaf error has norm at most
$\varepsilon_D$,
\[
 \norm{E_D}_2\le\varepsilon_D,\qquad
 \norm{E_\ell}_2\le\varepsilon_\ell,\qquad
 \norm E_2\le\varepsilon_D+\sum_{\ell=0}^{L-1}\varepsilon_\ell.
\]
The least-common-ancestor partition is also a disjoint partition of
matrix entries. Summing $|E_{ij}|^2$ over that partition proves the exact
identity
\[
 \norm E_F^2=
 \sum_{\rm leaf}\norm{E[I_\tau,I_\tau]}_F^2
 +\sum_{\rm sibling\ directions}\norm{E[I_\alpha,I_\beta]}_F^2.
\]
These are certificates for errors in the reconstructed physical blocks.
Errors in small coupling matrices alone require separate bounds on their
row and column bases.

\subsection{Nested HSS bases: two induction invariants and an entry count}

Use the data and ranks in the card: leaves have dense blocks $D_\tau$
and bases $U_\tau,V_\tau$ with $r_\tau\le k$ columns; for an internal
node with children $\alpha,\beta$,
\[
 U_\tau=\diag(U_\alpha,U_\beta)E_\tau,
 \qquad V_\tau=\diag(V_\alpha,V_\beta)F_\tau.
\]
The sibling blocks are
$U_\alpha B_{\alpha\beta}V_\beta^*$ and
$U_\beta B_{\beta\alpha}V_\alpha^*$. Put the root rank equal to zero.
The upward recursion computes
\[
 \widehat x_\tau=
 \begin{cases}
 V_\tau^*x_{I_\tau},&\tau\text{ a leaf},\\
 F_\tau^*\begin{bmatrix}\widehat x_\alpha\\\widehat x_\beta\end{bmatrix},
       &\tau\text{ internal}.
 \end{cases}
\]
Its induction invariant is exactly
$\widehat x_\tau=V_\tau^*x_{I_\tau}$, obtained by taking adjoints of
the nested-basis identity.

Start with a zero coefficient at the root. If $\widehat y_\tau$ has
already been accumulated, perform the downward assignments
\[
 \begin{bmatrix}\widehat y_\alpha\\\widehat y_\beta\end{bmatrix}
 =E_\tau\widehat y_\tau+
 \begin{bmatrix}
 B_{\alpha\beta}\widehat x_\beta\\
 B_{\beta\alpha}\widehat x_\alpha
 \end{bmatrix}.
\]
The invariant is
$U_\tau\widehat y_\tau=A[I_\tau,I_\tau^c]x_{I_\tau^c}$.
At the root both sides vanish. For a child, its complement splits into
the parent's complement and its sibling; the nested-basis identity
represents the first contribution and the sibling factorization the
second. Thus the invariant passes down the tree. On a leaf return
\[
 y_{I_\tau}=D_\tau x_{I_\tau}+U_\tau\widehat y_\tau.
\]
The invariant proves $y=Ax$.

With $q$ leaves, the exact parameter bound is
\[
 S=\sum_{\tau\ {\rm leaf}}(n_\tau^2+2n_\tau r_\tau)
 +\sum_{\tau\ {\rm internal}}
 \left[2(r_\alpha+r_\beta)r_\tau+2r_\alpha r_\beta\right].
\]
Every stored matrix is applied once in the two passes, so the same
accumulating matrix-vector count gives at most $2S$ operations. A full
binary tree with $q$ leaves has $q-1$ internal nodes. Therefore
\[
 S\le Nb+2Nk+6k^2(q-1).
\]
If $k\ge1$ and $k\le n_\tau\le2k$ at every leaf, then $qk\le N$,
and the sharper displayed consequence in the card follows:
$S\le2Nk+2Nk+6Nk=10Nk$. Root transfers have zero columns and contribute
no storage. Separate low ranks of sibling blocks do not by themselves
give the common nested bases assumed here.

\subsection{Compressed inversion: a projection removes the need for $C^{-1}$}

Let $D$ be invertible, $U,V\in\mathbb K^{N\times K}$, and
$M=V^*D^{-1}U$ invertible, with $0\le K\le N$. Put
\[
 R=D^{-1}U,\quad W=V^*,\quad H=M^{-1},\quad S=C+H,
 \qquad Q=RHW.
\]
Then $WR=M$, so $Q^2=Q$, $\operatorname{range}Q=\operatorname{range}R$,
and $\ker Q=\ker W$. In particular $R$ is injective, $W$ is surjective,
and
\[
 \mathbb K^N=\operatorname{range}R\oplus\ker W,\qquad\rank Q=K.
\]
The factorization $A=D+UCV^*=D(I+RCW)$ preserves this direct sum after
removing the invertible left factor $D$. On $\ker W$, $I+RCW$ is the
identity, while
\[
 (I+RCW)R=R(I+CM)=RSM.
\]
Consequently $A$ is invertible if and only if $S$ is invertible.

Define, exactly as in the card,
\[
 E=RH,\qquad F^*=HWD^{-1},\qquad
 G=(I-Q)D^{-1}=D^{-1}-D^{-1}UHV^*D^{-1}.
\]
The two short products
$AE=US$ and $AG=I-UF^*$ prove
\[
 A(ES^{-1}F^*+G)=I,
 \qquad A^{-1}=ES^{-1}F^*+G.
\]
Since right multiplication by $D^{-1}$ preserves rank and $Q$ is a
rank-$K$ projection, $\rank G=N-K$. If $D,U,V$ are conformally block
diagonal, so are $M,H,E,F,G$. For different block indices $i\ne j$,
\[
 (A^{-1})_{ij}=E_i(S^{-1})_{ij}F_j^*.
\]
The argument includes $K=0$: $S$ is the invertible empty matrix,
$ES^{-1}F^*=0$, and $G=D^{-1}$. At $K=N$, $Q=I$ and $G=0$.
Neither invertibility of $C$ nor nonsingularity of unspecified future
pivots is asserted. Invertibility of $A$ alone does not imply either of
the two input invertibility hypotheses.

\subsection{Inverse cut ranks: the nullspace maps are explicit}

Let $A$ be invertible, $B=A^{-1}$, and partition coordinates as $I\sqcup J$.
Both $A[I,J]$ and $B[I,J]$ have domain $\mathbb K^J$. There is a linear
bijection
\[
 \ker B[I,J]\longrightarrow\ker A[I,J],\qquad
 v\longmapsto B[J,J]v,
\]
with inverse $w\mapsto A[J,J]w$. Indeed, if $B[I,J]v=0$, then
$B(0,v)=(0,w)$, so applying $A$ gives
$A(0,w)=(0,v)$; this is exactly
$A[I,J]w=0$ and $A[J,J]w=v$. The converse is the same argument with
$A$ and $B$ exchanged. Therefore the nullities, hence the ranks, agree:
\[
 \rank B[I,J]=\rank A[I,J].
\]
Swap $I,J$ to obtain the other directed cut. No diagonal principal block
is required to be invertible. The proved declaration
\code{NLAlib.complementary_nullities_properties} already gives this
nullity result over $\R$; the displayed maps isolate its elementary
field-general extension.

For completeness, whole-cut ranks also produce nested bases. Define
$\mathcal U_\tau=\operatorname{range}B[I_\tau,I_\tau^c]$ and
$\mathcal V_\tau=\operatorname{range}(B[I_\tau^c,I_\tau])^*$.
If $\alpha$ is a child of $\tau$, restricting vectors in
$\mathcal U_\tau$ to $I_\alpha$ lands in $\mathcal U_\alpha$, because
$I_\tau^c\subseteq I_\alpha^c$. The analogous statement holds for
$\mathcal V$. Choose bases of these finite spaces. Expressing the
restrictions of each parent basis in its two child bases gives transfer
matrices and the required nesting. A sibling block has its columns in
$\mathcal U_\alpha$ and its adjoint columns in $\mathcal V_\beta$;
left inverses for the basis matrices therefore give a coupling matrix
$B_{\alpha\beta}$ and the exact factorization
$U_\alpha B_{\alpha\beta}V_\beta^*$. Row and column ranks may be
different; use separate ranks, or pad the smaller basis with zero columns
to their maximum, which is still at most $k$.
Applying the equal-cut-rank theorem to all clusters proves HSS/HBS
closure under inversion with the same rank bounds. This proves
existence of an exact representation, not a conditioning theorem or a
claim about an arbitrary inversion algorithm's cost.

\section{Effective resistance as a projector diagonal}
\label{sec:ns-resistance}

This section covers \code{laplacian-resistance-leverage}. Let $B$ be
the incidence matrix of a finite undirected loop-free graph, with any
orientation and positive weights $w_e$. Put
$W=\diag(w_e)$, $C=W^{1/2}B$, and $L=C^\top C$. A row of $B$ is
$b_e^\top$, so define $R_e=b_e^\top L^+b_e$.

The only matrix identity needed is
\[
 \Pi=CL^+C^\top=CC^+.
\]
For example, the Gram pseudoinverse identity gives
$L^+=C^+(C^+)^\top$, so the left side is
$(CC^+)(CC^+)^\top=CC^+$ by the Penrose identities. Thus $\Pi$ is
the orthogonal projector onto $\operatorname{range}C$. At coordinate $e$,
\[
 \Pi_{ee}=w_e b_e^\top L^+b_e=w_eR_e,
 \qquad
 \Pi_{ee}=\norm{\Pi e_e}_2^2\in[0,1].
\]
The upper bound uses the orthogonal decomposition
$e_e=\Pi e_e+(I-\Pi)e_e$. Taking the trace gives
\[
 \sum_e w_eR_e=\tr\Pi=\rank C=\rank L.
\]
Finally,
\[
 x^\top Lx=\sum_e w_e(b_e^\top x)^2
\]
shows that $\ker L$ consists exactly of vectors constant along each
edge, hence constant on each connected component. Conversely those
vectors are visibly in the kernel. The indicator vectors of the $c$
components are independent and span this kernel, including isolated
vertices. Rank--nullity yields $\rank L=|V|-c$ and Foster's identity.

Changing orientations multiplies the corresponding rows of $C$ by
$-1$, leaving $L$ and all projector diagonal entries unchanged. This
argument includes disconnected graphs, edgeless graphs, and the empty
graph. The resistance formula here concerns endpoints joined by an edge;
it does not identify a finite cross-component pseudoinverse quadratic
form with electrical resistance across disconnected components.

\paragraph{Verified reuse.}
The snapshot contains
\code{NLAlib.moorePenroseInverse_conjTranspose_mul_self},
\code{NLAlib.rangeProjector_eq_mul_moorePenroseInverse},
\code{NLAlib.rangeProjector_isHermitian},
\code{NLAlib.rangeProjector_isIdempotentElem}, and
\code{NLAlib.trace_mul_moorePenroseInverse}.
Only the incidence-kernel/component bridge is graph-specific; all
resistance bounds are finite projector algebra.

\section{Normal-matrix mixed moments as a finite cross-Gram identity}
\label{sec:ns-mixedmoments}

This covers \code{normal-matrix-mixed-chebyshev-moments}. Let
$A\in\C^{n\times n}$, $n\ge1$, be normal, with $\norm A_2\le1$,
and $\norm b_2=1$. Put $H=(A+A^*)/2$ and $K=(A-A^*)/(2i)$.
Choose a unitary diagonalization
$A=U\diag(\lambda_s)U^*$, supplied by the proved declaration
\code{NLAlib.exists_unitary_diagonalization_of_normal}. Then
\[
 H=U\diag(\operatorname{Re}\lambda_s)U^*,\qquad
 K=U\diag(\operatorname{Im}\lambda_s)U^*.
\]
Thus $H,K$ are commuting Hermitian contractions. For every distinct
eigenvalue $\lambda$, let $P_\lambda$ be its orthogonal eigenspace
projector. The finite probability measure
\[
 \mu_{A,b}=\sum_{\lambda\in\sigma(A)}
       \norm{P_\lambda b}_2^2
       \delta_{(\operatorname{Re}\lambda,\operatorname{Im}\lambda)}
\]
is manifestly independent of the eigenbasis. Resolving each projector
into any orthonormal basis recovers the card's sum
$\sum_s|u_s^*b|^2\delta_{(\operatorname{Re}\lambda_s,
\operatorname{Im}\lambda_s)}$. Orthogonal eigenspaces exhaust the
space, so the weights sum to $\norm b_2^2=1$. All support points lie in
the closed unit disk. This handles repeated eigenvalues without selecting
one basis canonically.

For arbitrary complex-coefficient polynomials $p,q$, polynomial
evaluation of the two diagonalizations gives the finite identity
\[
 \int\overline{p(x)}q(y)\,d\mu_{A,b}
 =\sum_s |u_s^*b|^2\,
     \overline{p(\operatorname{Re}\lambda_s)}
                  q(\operatorname{Im}\lambda_s)
 =b^*p(H)^*q(K)b.
\]
Thus no measure-theoretic limiting argument is required. Define
$\widetilde T_0=1/\sqrt\pi$ and
$\widetilde T_j=\sqrt{2/\pi}\,T_j$ for $j\ge1$; put
$X_j=\widetilde T_j(H)b$ and $Y_j=\widetilde T_j(K)b$.
For $0\le j,k\le M$ this gives
\[
 \Gamma_{jk}=\int\widetilde T_j(x)\widetilde T_k(y)\,d\mu_{A,b}
            =X_j^*Y_k,\qquad \Gamma=X^*Y.
\]

The normalized recurrence needs its special initial step:
\[
 X_0=b/\sqrt\pi,\qquad X_1=\sqrt{2/\pi}\,Hb,
 \qquad X_2=2HX_1-\sqrt2\,X_0,
\]
and $X_{j+1}=2HX_j-X_{j-1}$ for $j\ge2$; use the same formulas with
$K$ for $Y$. Omit unavailable indices when $M=0$ or $M=1$.
Each of the two sequences needs $M$ applications of its Hermitian
operator. Computing $Hv=(Av+A^*v)/2$ or
$Kv=(Av-A^*v)/(2i)$ requires two queries, giving at most $4M$ queries to
$A$ and $A^*$. The two initial actions on $b$ can share queries, but the
advertised upper bound already holds. Recurrences use $O(nM)$ additional
arithmetic, and the $(M+1)^2$ finite inner products use
$O(n(M+1)^2)$. Retaining both vector sequences and all output moments
uses $O(n(M+1)+(M+1)^2)$ scalars. A smaller storage claim requires streamed
output or an additional relation between $M$ and $n$.



\section{Data fits and reduced models as finite factorizations}
\subsection{DMD consistency is a kernel inclusion}
\label{route:dmd-linear-consistency}
\textbf{Card:} \code{dmd-linear-consistency}.
For $X,Y\in\C^{n\times m}$, set $A=YX^\dagger$ and $P=X^\dagger X$.
The Moore--Penrose identities say $P$ is the orthogonal projector onto
$(\ker X)^\perp$. Consequently
\[
 Y-AX=Y(I-P),\qquad AX=Y\ \Longleftrightarrow\ \ker X\subseteq\ker Y.
\]
Indeed $I-P$ has range $\ker X$. Every exact fit $M$ agrees with $A$
on $\operatorname{range}X$, since $(M-A)X=0$. More generally, orthogonal
splitting on the sample-coordinate space gives
\[
 \norm{Y-MX}_F^2=\norm{Y(I-P)}_F^2+\norm{YP-MX}_F^2.
\]
Thus $A$ is a least-squares fit. Every minimizer is
$A+Z(I-XX^\dagger)$; its two summands are Frobenius orthogonal, so $A$
has minimum Frobenius norm. For consecutive snapshots
$X=[x_0,\ldots,x_{m-1}]$, $Y=[x_1,\ldots,x_m]$, consistency implies
$x_j=A^jx_0$ for $0\le j\le m$ by a finite induction.
It imposes no condition on unobserved future snapshots.

\subsection{Reduced DMD eigenpairs: use $BU^*$ and $U^*B$}
\label{route:dmd-reduced-eigenpair-lifting}
\textbf{Card:} \code{dmd-reduced-eigenpair-lifting}.
Retain $r$ positive singular values of $X$, with associated $U,V,\Sigma$.
Set $B=YV\Sigma^{-1}$, $A_r=BU^*$ and $\widetilde A=U^*B$.
Complete $U$ to a unitary basis $[U,U_\perp]$. In this basis
\[
 A_r\sim\begin{bmatrix}\widetilde A&0\\U_\perp^*B&0\end{bmatrix},\qquad
 \det(zI_n-A_r)=z^{n-r}\det(zI_r-\widetilde A).
\]
This proves the multiplicity statement without diagonalizability.
If $\widetilde Aw=\lambda w$ and $\lambda\ne0$, then $Bw\ne0$ and
$A_r(Bw)=\lambda Bw$. Conversely, $A_r\phi=\lambda\phi$, $\lambda\ne0$,
implies $U^*\phi\ne0$ and $\widetilde A(U^*\phi)=\lambda U^*\phi$.
The projected vector $Uw$ satisfies the exact residual identity
\[
 A_rUw-\lambda Uw=(I-UU^*)Bw.
\]
Zero eigenvalues need the determinant formula, not a claim that these
eigenvector maps remain injective. If $r=0$, both constructions are empty
and $A_r=0$.

\subsection{ERA realization: expose a single invertible change of state}
\label{route:hankel-svd-era-realization}
\textbf{Card:} \code{hankel-svd-era-realization}.
Let $\mathcal O_s=[C;CA;\ldots;CA^{s-1}]$ and
$\mathcal R_t=[B,AB,\ldots,A^{t-1}B]$ both have rank $n$.
Write
\[
 H_0=\mathcal O_s\mathcal R_t=U\Sigma V^*,\quad
 H_1=\mathcal O_sA\mathcal R_t,\quad
 T=\mathcal R_tV\Sigma^{-1/2}.
\]
Here the SVD is compact, with $n$ positive entries. Since
$\mathcal O_sT=U\Sigma^{1/2}$, $T$ is injective and square, hence invertible.
Also $T\Sigma^{1/2}V^*=\mathcal R_t$: multiply by the injective
$\mathcal O_s$ to check it. It follows directly that
\[
 \widehat A=\Sigma^{-1/2}U^*H_1V\Sigma^{-1/2}=T^{-1}AT.
\]
Let $E_q$ select the first input block of columns and $E_p$ the first output
block of rows. Then
\[
 \widehat B=\Sigma^{1/2}V^*E_q=T^{-1}B,\qquad
 \widehat C=E_p^*U\Sigma^{1/2}=CT.
\]
Thus $\widehat C\widehat A^j\widehat B=CA^jB$ for every $j\ge0$.
Any realization of these data has state dimension at least $\rank H_0=n$,
because its Hankel matrix factors through its state space. This proves
minimality as a rank bound, without developing realization equivalence
theory. Stability is not an assumption of this identity.

\subsection{Gramian balancing through an SVD of the factors}
\label{route:gramian-svd-balancing}
\textbf{Card:} \code{gramian-svd-balancing}.
Let $P=RR^*$, $Q=O^*O$, $H=OR$, and choose any $r$ positive singular
triplets, with orthonormal columns $U,V$, so that
$HV=U\Sigma$ and $H^*U=V\Sigma$.
The block need not retain all positive singular values. Define
\[
 \Psi=RV\Sigma^{-1/2},\qquad \Phi=O^*U\Sigma^{-1/2}.
\]
Multiplication of these factorizations gives
\[
 \Phi^*\Psi=I_r,\quad \Phi^*P\Phi=\Sigma,\quad
 \Psi^*Q\Psi=\Sigma. \tag{B1}
\]
For example $\Phi^*R=\Sigma^{1/2}V^*$ and $O\Psi=U\Sigma^{1/2}$,
which prove both Gram identities. If $P,Q\succ0$ and all $n$ state
directions are retained, so $r=n$,
$T=\Psi$ is invertible, $T^{-1}=\Phi^*$, and
\[
 T^{-1}PT^{-*}=\Sigma=T^*QT,\qquad T^{-1}PQT=\Sigma^2.
\]
In the semidefinite case, the nonzero eigenvalues of $PQ$ still equal the
positive squared singular values of $OR$, with algebraic multiplicities.
Indeed $R^*O^*OR=(OR)^*(OR)$, and the determinant identity for $UV$ and
$VU$ identifies its nonzero characteristic roots with those of
$RR^*O^*O=PQ$. Repeated singular values merely change the SVD bases.
When $r=0$, (B1) is an identity of empty matrices, not an invertible state
transformation.

\subsection{Balanced truncation: one square completion gives the factor $2$}
\label{route:balanced-truncation-energy-bound}
\textbf{Card:} \code{balanced-truncation-energy-bound}.
Assume a positive diagonal $\Sigma=\diag(\sigma_1,\ldots,\sigma_n)$ satisfies
\[
 A\Sigma+\Sigma A^*=-BB^*,\qquad
 A^*\Sigma+\Sigma A=-C^*C. \tag{B2}
\]
First discard a block for which $\Sigma_2=\tau I$, $\tau>0$.
Partition the state as $(x_1,x_2)$ and let the retained state $x_r$ satisfy
$\dot x_r=A_{11}x_r+B_1u$. All states start at zero. Put
\[
 x_-=(x_1-x_r,x_2),\quad x_+=(x_1+x_r,x_2),\quad
 v=A_{21}x_r+B_2u,\quad z=B^*\Sigma^{-1}x_+.
\]
The exact state equations are
$\dot x_-=Ax_-+(0,v)$ and $\dot x_+=Ax_++2Bu-(0,v)$.
The output error is $e=Cx_-$. The second identity of (B2) gives
\[
 \frac d{dt}(x_-^*\Sigma x_-)=-\norm e^2+2\tau\operatorname{Re}(x_2^*v).
\]
Congruencing the first identity of (B2) by $\Sigma^{-1}$ gives
$A^*\Sigma^{-1}+\Sigma^{-1}A=-\Sigma^{-1}BB^*\Sigma^{-1}$, hence
\[
 \frac d{dt}(x_+^*\Sigma^{-1}x_+)
 =-\norm z^2+4\operatorname{Re}(u^*z)
                  -\frac2\tau\operatorname{Re}(x_2^*v).
\]
Multiplying by $\tau^2$ and adding cancels the problematic coupling exactly:
\[
 \frac d{dt}\left(x_-^*\Sigma x_-+\tau^2x_+^*\Sigma^{-1}x_+\right)
 =-\norm e^2+4\tau^2\norm u^2-\tau^2\norm{z-2u}^2. \tag{B3}
\]
The energy starts at zero and is nonnegative. Integrating over $[0,T]$
proves $\norm e_{L^2(0,T)}\le2\tau\norm u_{L^2(0,T)}$.
For $u\in L^2(0,T)$, the convolution states are absolutely continuous;
the identities hold almost everywhere and the same fundamental theorem
of calculus applies. Thus smooth-input approximation is optional.

Now discard the distinct tail values $\tau$ one group at a time. Every
principal truncation inherits both equations (B2), because $\Sigma$ is
block diagonal. The triangle inequality proves
\[
 \boxed{\norm{y-y_r}_{L^2(0,T)}
 \le2\!\sum_{\tau\in\{\sigma_{r+1},\ldots,\sigma_n\}}\!\tau\,
                       \norm u_{L^2(0,T)}
 \le2\sum_{j>r}\sigma_j\norm u_{L^2(0,T)}.} \tag{B4}
\]
The set in the first sum contains each distinct value once. A cut through
an equal-value group is allowed: discard its selected tail part as one
$\tau I$ block. The cases $r=0$ and $r=n$ are covered.
Monotone convergence of the squared integrals extends (B4) to the half-line.
If both full and reduced state matrices are Hurwitz, the frequency-domain
bound follows by feeding $u(t)=e^{i\omega t}v$ on $[0,T]$, dividing by
$\sqrt T$, and letting $T\to\infty$: exponentially decaying initial
transients contribute $o(\sqrt T)$ and the steady output is the transfer
matrix at $i\omega$ times $v$. This avoids Plancherel as well.

\subsection{DEIM: a Schur complement and an oblique projector}
\label{route:deim-greedy-interpolation}
\textbf{Card:} \code{deim-greedy-interpolation}.
Let $U=[u_1,\ldots,u_r]$ have orthonormal columns. At step $j$, solve the
already invertible interpolation system on the previous selected indices,
and subtract that combination of $u_1,\ldots,u_{j-1}$ from $u_j$.
The residual $w_j$ vanishes at the old indices and is not the zero vector,
by linear independence of the columns. An index maximizing $|w_j(i)|$
is therefore new and has nonzero residual. The determinant of the enlarged
sample matrix is the old determinant times this residual, by one column
operation. Induction proves invertibility of $M=P^*U$ and distinctness of
all $r$ indices; no generic-position hypothesis is used.
For $\Pi=UM^{-1}P^*$, direct products prove
\[
 \Pi^2=\Pi,\qquad \Pi U=U,\qquad P^*\Pi=P^*,\qquad
 \norm\Pi_2=\norm{M^{-1}}_2.
\]
The norm equality uses that $U$ is an isometry and $P^*$ a coisometry.
For $0<r<n$, decompose orthogonally into $\operatorname{range}U$ and its
complement. Then
\[
 \Pi=\begin{bmatrix}I&K\\0&0\end{bmatrix},\quad
 I-\Pi=\begin{bmatrix}0&-K\\0&I\end{bmatrix},\quad
 \norm{I-\Pi}_2=\norm\Pi_2=\sqrt{1+\norm K_2^2}.
\]
Consequently
$\norm{f-\Pi f}_2\le\norm{M^{-1}}_2\norm{(I-UU^*)f}_2$.
For $r=n$ the error is zero. For $r=0$ the sample system is empty and
$\Pi=0$; use the trivial bound $\norm f\le\norm f$, rather than assigning
the inverse of an empty matrix a nonzero operator norm.

\section{Polynomial polar iteration without a composition theorem}
\label{route:polar-minimax-polynomial-iteration}
\textbf{Card:} \code{polar-minimax-polynomial-iteration}.
Here is a direct proof of the advertised error estimate that does not need
the stronger theorem that greedy composition is globally optimal.
Fix $q\ge1$ and define
\[
 c_j=4^{-j}\binom{2j}{j},\qquad
 h_q(x)=x\sum_{j=0}^qc_j(1-x^2)^j.
\]
The finite coefficient recurrence $c_{j+1}=(2j+1)c_j/(2j+2)$ gives
$h_q'(x)=(2q+1)c_q(1-x^2)^q$. Hence $0<h_q(x)\le1$ for $0<x\le1$.
For $0\le u\le1$, the coefficients of degree $0,\ldots,q$ in
$(\sum_{j=0}^qc_ju^j)^2$ are all $1$, while the remaining coefficients
are nonnegative. For a fully finite proof write $C(u)=\sum_{j=0}^qc_ju^j$
and $D(u)=C(u)^2=\sum_jd_ju^j$. The recurrence says
$2(1-u)C'(u)-C(u)$ has zero coefficients through degree $q-1$.
Multiplication by $C$ shows that $(1-u)D'(u)-D(u)$ has the same
vanishing coefficients. Thus
$(j+1)d_{j+1}=(j+1)d_j$ for $0\le j<q$; since $d_0=1$,
$d_0=\cdots=d_q=1$. No infinite binomial series is needed.
It follows that
\[
 h_q(x)^2\ge(1-u)(1+u+\cdots+u^q)=1-u^{q+1},\quad u=1-x^2. \tag{P1}
\]

On any positive interval $[a,b]$ with $a<b$, an odd degree-at-most-$2q+1$ minimax
polynomial exists by finite-dimensional compactness: values at $q+1$
distinct positive nodes determine its coefficients through a weighted
Vandermonde matrix, so a bounded error sublevel has bounded coefficients.
A linear competitor gives error less than one. Its optimal image
$[a',b']$ is therefore positive and satisfies $a'+b'=2$: otherwise
multiplying the polynomial by $2/(a'+b')$ strictly improves its error.
The optimal polynomial is nonconstant, since a constant odd polynomial
is zero and has error one. Therefore $a'<b'$, and every subsequent
interval again has positive length. A singleton input interval can
separately be sent exactly to one by a linear polynomial; it does not
occur starting from $0<\ell<1$.
Write $\rho=a/b$ and $\delta=1-\rho^2$. By (P1), $h_q(x/b)$ has image
$[m,1]$, with $m\ge\sqrt{1-\delta^{q+1}}$. Scaling this competitor by
$2/(1+m)$ shows that the minimax error $\epsilon$ satisfies
\[
 \epsilon\le\frac{1-m}{1+m},\qquad
 \frac{a'}{b'}=\frac{1-\epsilon}{1+\epsilon}\ge m.
\]
To fix indices, let $I_1=[\ell,1]$, $I_{t+1}=p_t(I_t)$,
$\rho_t=\min I_t/\max I_t$ and $\delta_t=1-\rho_t^2$.
Write $\ell_{t+1}=\min I_{t+1}$; its upper endpoint is
$2-\ell_{t+1}$. Thus successive greedy interval images satisfy
\[
 \delta_{t+1}\le\delta_t^{q+1},\qquad
 \epsilon_t=1-\ell_{t+1}
 =\frac{1-\rho_{t+1}}{1+\rho_{t+1}}\le\delta_{t+1}.
\]
Starting with $[\ell,1]$, this proves
$1-\ell_{T+1}\le(1-\ell^2)^{(q+1)^T}$.
Continuity ensures that each interval image is the full interval used at
the next step. This argument never assumes that unnormalized minimax
iterates have upper endpoint at most one.

For a compact SVD $A=U_r\Sigma_rV_r^*$ and any odd polynomial
$p(x)=\sum_{j=0}^qa_jx^{2j+1}$, orthonormal-column cancellation proves
\[
 \sum_{j=0}^qa_jA(A^*A)^j=U_rp(\Sigma_r)V_r^*.
\]
Induction therefore gives $X_T=U_rP_T(\Sigma_r)V_r^*$, even for signed
intermediate scalar values. For the canonical partial polar factor
$Q=U_rV_r^*$,
\[
 \norm{Q-X_T}_2=\max_{i\le r}|1-P_T(\sigma_i)|
 \le1-\ell_{T+1}\le(1-\ell^2)^{(q+1)^T}.
\]
Zero singular directions stay zero; for $A=0$ both sides of the matrix
error identity are zero. The interval supremum is attained on a rank-one
matrix with its nonzero singular value at a scalar maximizer, proving
the worst-case equality for any fixed positive dimensions.
For Horner evaluation, form $S=X^*X$ using one product, evaluate
$a_qS^q+\cdots+a_0I$ using $q-1$ products (the first multiplication by
the scalar matrix is free), and multiply by $X$ using one product.
The total is at most $(q+1)T$ matrix products. Coefficient computation and
floating-point error are separate questions. The exponent here is
$(q+1)^T$, not $(2q+2)^T$.


\section{Stahl's trace-exponential theorem: the remaining deep core}
\label{route:trace-exponential-stahl}
\textbf{Card:} \code{trace-exponential-stahl}.
\textbf{Status of this reroute: partial.} The consequences and an exact
finite-inequality interface are derived below. Positivity of that interface
for general noncommuting matrices is not proved here. It is equivalent to
the hard part of BMV; neither convexity of $\tr e^A$ nor the elementary
Fr\'echet derivative formula establishes it. This card is included, not
deferred or silently marked complete.

Let $A,H$ be Hermitian, $d\ge1$, $F(t)=\tr\exp(A+tH)$, and put
$a=\lambda_{\min}(H)$, $b=\lambda_{\max}(H)$. For a real polynomial
$p(x)=\sum_kp_kx^k$, define the finite linear functional
\[
 \mathcal L(p)=\sum_kp_kF^{(k)}(0).
\]
There is no measure in this definition. All its inputs are derivatives
of a finite matrix exponential. The iterated product rule, or its
ordered-simplex integral, implies
$|F^{(k)}(0)|\le d e^{\norm A_2}\norm H_2^k$; each product contains $k$
copies of $H$, and the ordered simplex has volume $1/k!$, cancelling the
derivative factorial. The scalar function $F$ is entire, by the normally
convergent exponential series of the affine matrix pencil.

\subsection{An exact family of finite inequalities}
For $a<b$, consider the following assertion, for all nonnegative integers
$j,k$:
\begin{equation}\label{eq:bmv-finite-core}
 \mathcal L\bigl((x-a)^j(b-x)^k\bigr)\ge0. \tag{S1}
\end{equation}
It is equivalent to existence of the desired positive measure on $[a,b]$.
One direction follows by integrating a nonnegative polynomial. For the
other, let $y=(x-a)/(b-a)$ and define positive finite atomic measures
\[
 \nu_N=\sum_{r=0}^N w_{N,r}\,\delta_{a+(b-a)r/N},\qquad
 w_{N,r}=\binom Nr\mathcal L\bigl(y^r(1-y)^{N-r}\bigr),\quad N\ge1.
\]
Every weight is nonnegative by (S1), and the binomial theorem gives
$\sum_rw_{N,r}=\mathcal L(1)=F(0)$. For any fixed polynomial $p$,
\[
 \int p\,d\nu_N=\mathcal L(B_Np),
\]
where $B_N$ is the Bernstein polynomial operator on $[a,b]$.
For fixed degree, its coefficients converge to those of $p$: on $[0,1]$
this follows from
\[
 B_N(y^k)=N^{-k}\sum_{j=0}^k S(k,j)(N)_j y^j,
\]
using the elementary expansion of powers into falling factorials.
Therefore $\mathcal L(B_Np)\to\mathcal L(p)$.
Compactness of positive measures of fixed mass on a compact real interval
(equivalently, Helly selection for their distribution functions) supplies
a weakly convergent subsequence with positive limit $\nu$.
Testing the continuous monomials gives
$\int x^k\,d\nu=F^{(k)}(0)$ for all $k$.
Uniform absolute convergence of the scalar exponential series on $[a,b]$
then yields
\[
 \int e^{tx}\,d\nu(x)
 =\sum_{k=0}^\infty\frac{t^k}{k!}F^{(k)}(0)=F(t)
 \quad(t\in\R).
\]
This supplies a very narrow possible formal interface: prove the scalar
finite inequalities (S1), then reuse real measure compactness and finite
Bernstein identities. It does \emph{not} make (S1) elementary. In
particular, expanding powers into noncommuting words and declaring each
word trace nonnegative is invalid; positivity of their required sums is
precisely the issue.

\subsection{All advertised consequences after the positivity theorem}
The standard BMV support theorem for $B=bI-H\succeq0$ gives a measure
$\mu$ on $[0,b-a]$ representing $\tr e^{A-tB}$. Push it forward by
$s\mapsto b-s$. Since $bI$ commutes with every matrix,
$F(t)=e^{bt}\tr e^{A-tB}$, giving the required support $[a,b]$.
If the imported theorem is stated only for $t\ge0$, both sides extend
to entire functions because the support is compact, so the identity
holds for all real $t$ as well. Evaluation at zero gives
$\nu(\R)=\tr e^A$.
For every $k\ge0$, differentiating under the integral is justified locally
by the finite dominating constant
$\max(|a|,|b|)^k e^{(|t|+1)\max(|a|,|b|)}$, times the finite mass.
It gives
\[
 F^{(k)}(t)=\int x^ke^{tx}\,d\nu(x).
\]
For $r\ge0$, $t\ge0$, multiply the pointwise inequality
$e^{ta}\le e^{tx}\le e^{tb}$ by $x^{2r}\ge0$ and integrate:
\[
 0\le e^{ta}F^{(2r)}(0)\le F^{(2r)}(t)
                  \le e^{tb}F^{(2r)}(0).
\]
These consequences require no noncommutative differentiation after the
representation is available.

\subsection{Cases that genuinely close with finite linear algebra}
If $AH=HA$, simultaneous diagonalization gives
$F(t)=\sum_i e^{\alpha_i}e^{t\eta_i}$ and
$\nu=\sum_i e^{\alpha_i}\delta_{\eta_i}$.
For $a=b$ one has $H=aI$ and
$\nu=(\tr e^A)\delta_a$, without diagonalizing $A$.
Repeated eigenvalues cause no difficulty in these cases or in the moment
interface above. For dimension zero, if desired, use the zero measure
and avoid defining extreme eigenvalues.
For a fully general proof, Stahl's theorem or Eremenko's shorter exposition
still requires the positivity argument on the algebraic surface of the
characteristic pencil. This document does not replace that argument with
an unsupported finite-LA claim.

\section*{Sources and formalization handoff}
\addcontentsline{toc}{section}{Sources and formalization handoff}
The exact target hypotheses and source labels are those of the dated
\href{https://github.com/tchen-research/NLAlib/blob/6886918346afb81ba77e71fbbe5daa67c93d2bae/atlas/atlas.json}
{catalogue snapshot}. The proofs above are written out so that a source
citation is not used in place of the mathematical step being proposed.
Useful primary references for the harder groups are:
\begin{enumerate}
\item B. Beckermann and A. Townsend,
\href{https://arxiv.org/abs/1609.09494}{On the singular values of matrices
with displacement structure}. The bivariate identity (D1) isolates the
rank mechanism without requiring explicit Zolotarev extremizers.
\item V. Simoncini,
\href{https://www.dm.unibo.it/~simoncin/matrixeq.pdf}{Computational methods
for linear matrix equations}. The Gramian and residual identities above
separate finite algebra from semigroup or power convergence.
\item P. Benner and P. K\"urschner,
\href{https://csc.mpi-magdeburg.mpg.de/preprints/2013/MPIMD13-05.pdf}
{Computing real low-rank solutions of Sylvester equations by the factored
ADI method}. The shift convention here is fixed by the two displayed
half-step equations, so sign conventions can be checked directly.
\item M. Redmann,
\href{https://arxiv.org/abs/1907.05478}{An $L^2_T$-error bound for time-limited
balanced truncation}. The unrestricted Gramian specialization admits the
explicit square-completion identity (B3).
\item N. Amsel, D. Persson, C. Musco and R. M. Gower,
\href{https://arxiv.org/abs/2505.16932v5}{The Polar Express}.
The direct interval-ratio argument above proves the matrix card's rate
without importing global optimality of greedy compositions.
\item H. R. Stahl,
\href{https://arxiv.org/abs/1107.4875v3}{Proof of the BMV conjecture},
and A. Eremenko,
\href{https://arxiv.org/abs/1312.6003}{Herbert Stahl's proof of the BMV conjecture}.
These supply the deep positivity theorem left explicit in the last section.
\end{enumerate}

For implementation, first package the finite helper identities stated in
each section, then express the target card using the already proved
declarations listed there. No new axiom, weakened target, or undocumented
positivity assumption is proposed. In the single partial card, (S1) must
remain a named unproved prerequisite until a proof of BMV or that equivalent
core is available; the measure construction and derivative consequences
must not be mistaken for such a proof.
\end{document}
